\paragraph{Origin and Invariant Cancellation of the Normalization Constant: Endpoint Geometric Clarification of ``\texorpdfstring{$4$}{4}''}
This paragraph attributes the constant \(4\) that repeatedly appears in the disk compactification and endpoint scaling semigroup to conformal normalization, and emphasizes that it strictly cancels in observable invariants: \(\mathcal{B}_{-1}\) and the endpoint precision potential fix the ``time fiber/scale time'' as a one-dimensional additive ledger, while the constant \(4\) only changes the additive origin and carries no new dynamical degree of freedom.

\begin{proposition}[Origin of the Coefficient \(4\): The Square of the Conformal Factor \(2\) from the Cayley Normalization]\label{con:xi-four-from-cayley-conformal-factor}
Under the Cayley compactification \(w=\mathcal{C}(t)\) (\(t=x+\mathrm{i}y\in\mathbb{H}\)) fixed throughout this paper, we always have
\[
1+w=\frac{2}{1-\mathrm{i}t},
\qquad
|1+w|^{2}=\frac{4}{|1-\mathrm{i}t|^{2}}=\frac{4}{x^{2}+(1+y)^{2}},
\]
and from
\(
|w|^{2}=\frac{|1+\mathrm{i}t|^{2}}{|1-\mathrm{i}t|^{2}}
\)
we obtain
\[
1-|w|^{2}
=\frac{4y}{x^{2}+(1+y)^{2}}.
\]
Therefore, the coefficient \(4\) appearing in \(|1+w|^{2}\), \(1-|w|^{2}\), and the comoving depth closed form
\(
h_{x_{0}}(\rho)=\frac{4\delta}{(\gamma-x_{0})^{2}+(1+\delta)^{2}}
\)
is not a new invariant, but rather the square of the conformal factor \(2\) of the Cayley transform under this normalization (the proof of Proposition \ref{prop:app-busemann-poisson-minus1} and Theorem \ref{thm:app-comoving-horizon-localization}).
\end{proposition}

\begin{proposition}[Renormalizability of ``\(4\)'': Under a More General Cayley Gauge, the Constant Rescales by \texorpdfstring{$\lambda^{2}$}{lambda^2}]\label{con:xi-four-rescales-under-cayley-gauge}
Let \(\lambda>0\), and consider the equivalent Cayley family that preserves the endpoint reference but changes the linear gauge:
\[
\mathcal{C}_{\lambda}(t):=\mathcal{C}(t/\lambda),
\qquad
w=\mathcal{C}_{\lambda}(t),\quad t\in\mathbb{H}.
\]
Then
\[
1+w=\frac{2\lambda}{\lambda-\mathrm{i}t},
\qquad
|1+w|^{2}=\frac{4\lambda^{2}}{|\lambda-\mathrm{i}t|^{2}},
\]
so that the \(4\) appearing in the same derivation uniformly rescales to \(4\lambda^{2}\) (while the numerator of \(1-|w|^{2}\) scales linearly in \(\lambda\)).
Therefore, the constant \(4\) is a coordinate normalization constant; genuine geometric quantities can only appear in ratios and logarithmic differences that are invariant under such rescalings (see Proposition \ref{con:xi-busemann-ratio-cancels-four} and Proposition \ref{con:xi-four-only-shifts-additive-origin}).
\end{proposition}

\begin{proposition}[Location of Invariants: \(4\) Strictly Cancels in Busemann/Poisson Ratios]\label{con:xi-busemann-ratio-cancels-four}
Define the exponentiated Busemann/Poisson coordinate at endpoint \(-1\) (Proposition \ref{prop:app-busemann-poisson-minus1})
\[
\mathcal{B}_{-1}(w)=\frac{1-|w|^{2}}{|1+w|^{2}}.
\]
In the upper half-plane coordinates \(t=\mathcal{C}^{-1}(w)=x+\mathrm{i}y\), dividing the two identities containing \(4\) from Proposition \ref{con:xi-four-from-cayley-conformal-factor} yields the strict invariant
\[
\boxed{\;\mathcal{B}_{-1}(w)=y=\Im(t).\;}
\]
Therefore, the constant \(4\) completely cancels in visible invariants; its sole role is to ensure that \(\mathcal{B}_{-1}\) is exactly identified with the half-plane height \(y\), thereby fixing the ``off-slice displacement/time fiber coordinate'' as the intrinsic readout of endpoint weight.
\end{proposition}

\begin{proposition}[Formal Similarity with the ``Area Law'': The \texorpdfstring{$2^{2}$}{2^2} Factor of the Same Class of Conformal Projections]\label{con:xi-four-as-universal_conformal_square}
In higher-dimensional compactification/projective coordinates, the same conformal factor \(\Omega(x)=\frac{2}{1+|x|^{2}}\) and its squares or powers also appear; hence the pullback densities of spherical metrics and area elements all carry ``\(4\)''-type constants.
For example, in the two-dimensional case, stereographic projection yields the standard identity
\[
\dd s^{2}=\frac{4}{(1+|z|^{2})^{2}}\,|\dd z|^{2},
\qquad
\dd A=\frac{4}{(1+|z|^{2})^{2}}\,\dd A_{\RR^{2}}.
\]
This phenomenon indicates that the \(4\) appearing in the endpoint compactification of this paper is a universal square constant of conformal normalization, not an invariant specific to any particular arithmetic object.
\end{proposition}

\begin{proposition}[Another Occurrence of ``\(4\)'' in the Endpoint Scaling Semigroup: The M\"obius Discriminant Locks the Scale Time]\label{con:xi-four_in_mobius_discriminant}
The M\"obius automorphism \(\psi_m\) induced by scale dilation (\eqref{eq:scale-psi-m}), when written as \(\psi_m(u)=(au+b)/(cu+d)\), satisfies
\(
a=d=m+1,\ b=c=1-m
\), so its discriminant satisfies
\[
ad-bc=(m+1)^{2}-(1-m)^{2}=4m
\]
(Proposition \ref{prop:scale-linearization-identity-multipliers}).
This ``\(4\)'' is likewise an algebraic image of the endpoint normalization, and is compatible with the endpoint precision potential
\(
y(u)=-\log|\mathfrak{w}(u)|
\)
and its strict additive time law
\(
y(\psi_m(u))=y(u)+\log m
\)
(\eqref{eq:endpoint-busemann-addition-law}): the matching between the scale parameter \(m\) and the endpoint time increment \(\log m\) is uniquely locked by this normalization at the group action level.
\end{proposition}

\begin{proposition}[Comparison with Black Hole Entropy: The Constant \(4\) Only Shifts the Additive Origin of the Entropy Potential]\label{con:xi-four-only-shifts-additive-origin}
Let \(b_{-1}(w):=\log\mathcal{B}_{-1}(w)\) be the Busemann function with basepoint at endpoint \(-1\) (up to a constant, Proposition \ref{prop:app-busemann-poisson-minus1}).
From Propositions \ref{con:xi-four-from-cayley-conformal-factor}--\ref{con:xi-busemann-ratio-cancels-four}, the constant \(4\) enters \(b_{-1}\) only through the constant term of \(\log|1+w|^{2}\), thereby changing only the additive origin without affecting any observable difference (such as scale time increments, spectral flow counts, defect flux, etc.).
Therefore, if ``entropy'' is understood as a logarithmic potential (additive ledger), then the role of the constant \(4\) is of the same type as a choice of units/normalization constant: its resemblance to the constant factor in the black hole entropy formula lies in both being fixed normalizations of the boundary--volume (or boundary--height) change of variables, rather than new dynamical degrees of freedom.
\end{proposition}

\begin{proposition}[Strict Version of ``Boundary Carries'': The Time Fiber Coordinate Is the Endpoint Poisson Weight]\label{con:xi-timefiber-equals-poisson-weight}
In the disk representation, the preimage of an equal-radius observation fiber \(|w|=\varrho\) in the upper half-plane simultaneously reaches \(y\downarrow 0\) and \(y\uparrow\infty\), and in the high-height regime forces contributions to concentrate in the endpoint neighborhood \(w\approx -1\) (endpoint compactification/boundary layer compression law; see Corollary \ref{cor:app-offcritical-horocycle-busemann} and Corollary \ref{cor:app-horocycle-visible-resolution}).
Furthermore, the time fiber coordinate satisfies the strict identity
\[
y=\Im(\mathcal{C}^{-1}(w))=\mathcal{B}_{-1}(w)=\frac{1-|w|^{2}}{|1+w|^{2}}
\]
(Proposition \ref{prop:app-busemann-poisson-minus1}).
Therefore, the mathematical content of ``boundary carrying/black-hole-type compression'' is: the visible continuous coordinate is not additionally introduced, but given by the Poisson/Busemann weight at endpoint \(-1\); the constant \(4\) appearing therein is entirely from conformal normalization and is strictly absorbed at the invariant level.
\end{proposition}

\begin{proposition}[Curvature Normalization Law: The Coefficient \(4\) Is Equivalent to the Unique Metric Scale for Constant Curvature \(-1\)]\label{con:xi-curvature-normalization-alpha4}
Take the hyperbolic metric on the upper half-plane model
\[
\dd s_{\mathbb{H}}^{2}=\frac{\dd x^{2}+\dd y^{2}}{y^{2}}
\]
(Gaussian curvature \(K\equiv -1\)).
For any M\"obius transformation \(w=w(t)\) conformally mapping \(\mathbb{H}\) to the unit disk \(\mathbb{D}\), the hyperbolic metric on the disk must be written as
\[
\dd s_{\mathbb{D}}^{2}=\frac{4\,|\dd w|^{2}}{(1-|w|^{2})^{2}},
\]
which also satisfies \(K\equiv -1\).
More generally, if rewritten as
\(
\dd s_{\mathbb{D}}^{2}=\frac{\alpha\,|\dd w|^{2}}{(1-|w|^{2})^{2}}
\)
(\(\alpha>0\) constant), then this metric is merely a constant scaling: \(\dd s_{\mathbb{D},\alpha}^{2}=\frac{\alpha}{4}\dd s_{\mathbb{D}}^{2}\), so the curvature transforms inversely by
\[
\boxed{\;K_{\alpha}=-\frac{4}{\alpha}\; }.
\]
Therefore, \(\alpha=4\) is the unique scale choice that ``normalizes the curvature to \(-1\)''; this condition is independent of any arithmetic object and is determined solely by fixing geometric units.
\end{proposition}

\begin{proposition}[Horizon Measure Normalization: Unit Mass of the Poisson Kernel and Compatible Locking with the Curvature Scale]\label{con:xi-poisson-kernel-unit-mass-normalization}
Define the Poisson kernel on the unit disk
\[
P_{w}(e^{\mathrm{i}\theta})
:=\frac{1-|w|^{2}}{|e^{\mathrm{i}\theta}-w|^{2}}
\qquad(w\in\mathbb{D}),
\]
then the strict normalization identity holds:
\[
\boxed{\;
\frac{1}{2\pi}\int_{0}^{2\pi}P_{w}(e^{\mathrm{i}\theta})\,\dd\theta=1\; }.
\]
This unit mass fixes the boundary horizon measure as a probability measure (harmonic measure), thereby determining, in parallel with the curvature normalization of Proposition \ref{con:xi-curvature-normalization-alpha4}, an implementation-independent set of ``horizon conventions'':
constant scaling would only introduce same-order cancellable unit transformations in the ``density \(\times\) measure'' decomposition, without entering any auditable invariant ratio.
\end{proposition}

\begin{proposition}[Coordinate Invariance of the Defect Entropy Functional: \(4\) Serves Only as a Unit Conversion Constant for Entropy]\label{con:xi-defect-entropy-coordinate_invariant}
Let point defects produce the profile under comoving scanning:
\[
H(x)=\sum_{j=1}^{\kappa} m_j\,\frac{4\delta_j}{(x-\gamma_j)^2+(1+\delta_j)^2}.
\]
Define the dimensionless defect entropy (also interpretable as the normalized area of defect flux)
\[
S_{\mathrm{def}}
:=\frac{1}{4\pi}\int_{\RR}H(x)\,\dd x.
\]
Then by \(\int_{\RR}\frac{\dd x}{(x-\gamma)^2+a^2}=\pi/a\), the strict closed form is obtained:
\[
\boxed{\;
S_{\mathrm{def}}
=\sum_{j=1}^{\kappa}m_j\,\frac{\delta_j}{1+\delta_j}\in(0,\kappa)\; }.
\]
Therefore, the constant \(4\) completely cancels in \(S_{\mathrm{def}}\), serving only as a unit conversion from ``geometric density'' to ``dimensionless entropy''; the auditable information is concentrated in the ratio \(\delta/(1+\delta)\) and the integer multiplicity \(m\).
Moreover, under any conformal reparametrization that only changes the coordinate scale (e.g., the Cayley gauge change of Proposition \ref{con:xi-four-rescales-under-cayley-gauge} and the corresponding change of variables), the density and the integral measure scale at the same order, so \(S_{\mathrm{def}}\) remains invariant as a coordinate-independent scalar.
\end{proposition}

\begin{proposition}[Index--Entropy Decomposition: Strict Separation of the Integer Defect Index and the Continuous Defect Entropy]\label{con:xi-index-entropy-strict-separation}
Let the upper half-plane inner function \(\mathcal{S}\) be defined as in Proposition \ref{con:xi-spectral-shift-phase-density}, and denote its boundary phase winding number (integer defect index)
\[
\kappa=\sum_{j=1}^{\kappa}m_j
=\frac{1}{2\pi}\int_{\RR}\frac{\dd}{\dd x}\arg \mathcal{S}(x)\,\dd x\in\NN.
\]
Then the continuous defect entropy satisfies the strict identity
\[
\boxed{\;
S_{\mathrm{def}}
=\kappa-\sum_{j=1}^{\kappa}\frac{m_j}{1+\delta_j}\;}
\]
(obtained by summing term by term from \( \frac{\delta}{1+\delta}=1-\frac{1}{1+\delta}\)).
Therefore, the defect information permitted by the regime naturally splits into two classes of orthogonal invariants:
\begin{enumerate}
  \item \textbf{Integer index \(\kappa\)}: given by the total winding number of the boundary phase flow, topologically non-continuously adjustable;
  \item \textbf{Continuous entropy potential \(S_{\mathrm{def}}\)}: the visible degradation of \(\kappa\), whose deviation is continuously controlled by the ``depth penalty'' terms \(\sum \frac{m}{1+\delta}\).
\end{enumerate}
\end{proposition}

\begin{proposition}[``Area Law'' Counterpart: \texorpdfstring{$S_{\mathrm{def}}$}{Sdef} Is the Normalized Area of the Weighted Spectral Shift Density]\label{con:xi-sdef-as-weighted-spectral-shift-area}
Let the spectral shift density be
\[
W(x):=\frac{1}{2\pi}\frac{\dd}{\dd x}\arg \mathcal{S}(x),
\qquad \int_{\RR}W(x)\,\dd x=\kappa,
\]
and let the single-atom phase density be
\[
W_j(x):=\frac{1}{2\pi}\frac{\dd}{\dd x}\arg\!\left(\frac{x-\gamma_j-\mathrm{i}(1+\delta_j)}{x-\gamma_j+\mathrm{i}(1+\delta_j)}\right)
=\frac{1+\delta_j}{\pi\bigl((x-\gamma_j)^2+(1+\delta_j)^2\bigr)}.
\]
Then for the weight \(\omega_j:=\frac{\delta_j}{1+\delta_j}\in(0,1)\), the strict decomposition holds:
\[
\boxed{\;
\frac{1}{4\pi}H(x)=\sum_{j=1}^{\kappa}\omega_j\,m_j\,W_j(x)\; }.
\]
Therefore
\[
\boxed{\;
S_{\mathrm{def}}
=\frac{1}{4\pi}\int_{\RR}H(x)\,\dd x
=\int_{\RR}\left(\sum_{j=1}^{\kappa}\omega_j\,m_j\,W_j(x)\right)\dd x\; }.
\]
In this sense, \(S_{\mathrm{def}}\) is a normalized area law: it integrates the absolutely continuous part of the boundary phase flow (horizon current), weighted, into a dimensionless entropy; the coefficient \(4\) is solely responsible for normalizing this boundary integral to \(S_{\mathrm{def}}\in(0,\kappa)\), without introducing additional dynamical degrees of freedom.
\end{proposition}

\begin{proposition}[Limiting Mechanism: The Two-End Limits of \texorpdfstring{$\delta$}{delta} Yield an ``Entropy Collapse/Saturation'' Two-Phase Diagram]\label{con:xi-sdef-collapse-saturation-phase}
For a single-atom defect, its contribution to \(S_{\mathrm{def}}\) is \(m\,\frac{\delta}{1+\delta}\).
Therefore, the two-end limiting laws hold:
\[
\boxed{\;
\delta\downarrow 0\ \Longrightarrow\ S_{\mathrm{def}}\downarrow 0,
\qquad
\delta\uparrow\infty\ \Longrightarrow\ S_{\mathrm{def}}\uparrow \kappa\;}
\]
(the second being a formal parameter limit, characterizing the formal saturation ``far from the endpoint layer'').
This two-end limit provides a strict phase diagram for the endpoint compression semantics: when defects are pushed into the endpoint thin layer (\(\delta\to 0\)), the visible entropy flux collapses to \(0\); in the formal limit \(\delta\to\infty\), the visible entropy flux saturates at the integer index \(\kappa\).
The constant \(4\) in this phase diagram only fixes the entropy potential scale, aligning the saturation value strictly with the topological count.
\end{proposition}

\begin{proposition}[``Horizon Regularity'' Interpretation of the Constant \(4\): Period Identification and Cone-Defect-Free Normalization]\label{con:xi-four_as_horizon_regularization_no_cone}
In the spectral crossing formulation of the unitary-slice, the critical line parameter can be implemented as a boundary angular variable with \(2\pi\)-period identification (phase winding number is an integer).
When aligning this period identification with the half-plane linear coordinate through the Cayley/completion interface, the requirement that ``a complete traversal \(\Longleftrightarrow\) one period on the boundary'' fixes the conversion constant between angular and linear coordinates to a unique normalization; this normalization is compatible with the curvature \(-1\) scale of Proposition \ref{con:xi-curvature-normalization-alpha4} and the horizon measure convention of Proposition \ref{con:xi-poisson-kernel-unit-mass-normalization}, and ultimately manifests as the same square constant \(4\) in all visible formulas of Poisson/Busemann/scanning kernels.
Therefore, this constant can also be equivalently understood as a horizon regularity condition: it prevents the period identification from introducing additional cone defects (thus avoiding spurious indices/spurious entropy terms), and strictly compresses visible anomalies to the two classes of invariants ``integer index \(\kappa\) \(\times\) continuous depth entropy potential \(S_{\mathrm{def}}\).''
\end{proposition}

\paragraph{Residue Scale, Depth Weight, and Boundary Action: Another Set of Uniqueness Characterizations for ``\texorpdfstring{$4$}{4}''}

\begin{proposition}[Rigidity Characterization of the Coefficient \(4\): Unification of Spectral Residues and Geometric Depth Uniquely Determines the Constant]
Suppose point defects are given by \(\rho_j=\tfrac12-\delta_j+\mathrm{i}\gamma_j\) (\(\delta_j>0\)), and introduce on comoving coordinates a family of profile densities with an undetermined normalization (\(c>0\)):
\[
H_{c}(x):=\sum_{j} m_j\,\frac{c\,\delta_j}{(x-\gamma_j)^2+(1+\delta_j)^2}.
\]
Let its Cauchy transform (\(\Im z>0\))
\[
M_{H_c}(z):=\int_{\RR}\frac{H_c(x)}{x-z}\,\dd x,
\]
and denote the lower half-plane poles
\(
z_j^{\downarrow}:=\gamma_j-\mathrm{i}(1+\delta_j)
 \)
(corresponding to the Poisson/Busemann normalization of Proposition \ref{con:xi-timefiber-equals-poisson-weight} and the same analytic convention in Theorem \ref{thm:xi-offcritical-poisson-herglotz-inversion}).
Then by the fundamental integral identity
\(
\int_{\RR}\frac{1}{(x-\gamma)^2+a^2}\frac{1}{x-z}\,\dd x=\frac{\pi}{a}\frac{1}{\gamma-\mathrm{i}a-z}
 \),
we obtain the residue closed form
\[
\operatorname*{Res}_{z=z_j^{\downarrow}}M_{H_c}(z)
=-\pi c\,m_j\,\frac{\delta_j}{1+\delta_j}.
\]
On the other hand, let the upper half-plane inner function
\[
\mathcal{S}(z):=\prod_{j}\left(\frac{z-\gamma_j-\mathrm{i}(1+\delta_j)}{z-\gamma_j+\mathrm{i}(1+\delta_j)}\right)^{m_j},
\qquad
M_{P}(z):=2\mathrm{i}\,\frac{\dd}{\dd z}\log\mathcal{S}(z).
\]
Then \(M_{P}\) has the residue at the same pole \(z_j^{\downarrow}\):
\[
\operatorname*{Res}_{z=z_j^{\downarrow}}M_{P}(z)=-2\mathrm{i}\,m_j.
\]
If we require that the ``dimensionless depth weight recovered from the spectral residue ratio'' coincides exactly with the geometric Busemann depth weight \(p(\delta):=\delta/(1+\delta)\) (Proposition \ref{con:xi-defect-entropy-coordinate_invariant}), i.e., for all \(j\):
\[
\frac{\delta_j}{1+\delta_j}
=\frac{1}{2\pi}\left|\frac{\operatorname*{Res}_{z=z_j^{\downarrow}}M_{H_c}(z)}{\operatorname*{Res}_{z=z_j^{\downarrow}}M_{P}(z)}\right|,
\]
then necessarily and uniquely \(c=4\).
Therefore, the coefficient \(4\) is not an additional invariant, but the unique scale constant under which ``the spectral flow (given by the logarithmic derivative of \(\mathcal{S}\)) and the depth readout (given by the Cauchy transform of \(H_c\)) are exchangeable under the same regime normalization.''
\end{proposition}

\begin{proposition}[Uniqueness of Depth Weight: The Multiplicative Scaling Semigroup's M\"obius Linearization Forces \texorpdfstring{$p(\delta)=\delta/(1+\delta)$}{p(delta)=delta/(1+delta)}]
Consider the depth scaling action \(\delta\mapsto m\delta\) (\(m>0\)).
Suppose there exists a continuous monotone coordinate \(p:(0,\infty)\to(0,1)\) such that this scaling, in \(p\)-coordinates, appears for all \(m\) as a fractional linear transformation preserving the endpoints \(\{0,1\}\):
\[
p\longmapsto \Phi_m(p):=\frac{m p}{1+(m-1)p}.
\]
Then there must exist a constant \(a>0\) such that
\[
\boxed{\;
p(\delta)=\frac{\delta}{\delta+a}\;}
\]
(converting \(q(\delta):=\frac{p(\delta)}{1-p(\delta)}\) to a positive real axis coordinate yields \(q(m\delta)=m\,q(\delta)\), whence \(q(\delta)\) is linear in \(\delta\)).
Furthermore, under the unitary-slice normalization of this paper (fixing this free scale to \(a=1\)), we obtain the unique form
\[
\boxed{\;
p(\delta)=\frac{\delta}{1+\delta}\;}
\]
which, together with the above residue unification (thus uniquely locking \(c=4\)), jointly provides the coupled uniqueness of the ``visible entropy/visible flux'' weight and the coefficient \(4\).
\end{proposition}

\begin{proposition}[Boundary Action Formulation: Defect Entropy Is the Weighted Neumann Boundary Term of the Green Potential, with the Normalization Constant Uniquely Locked]
Denote \(t_j:=\gamma_j+\mathrm{i}(1+\delta_j)\in\mathbb{H}\), and define the Green-type potential
\[
U_j(t):=\log\left|\frac{t-\overline{t_j}}{t-t_j}\right|
\qquad(t\in\mathbb{H}).
\]
Then its outward normal derivative at the boundary \(y=0\) satisfies
\[
\partial_y U_j(x)=\frac{2(1+\delta_j)}{(x-\gamma_j)^2+(1+\delta_j)^2}.
\]
Define the weighted boundary action
\[
\mathcal{A}_{\mathrm{def}}
:=\frac{1}{4\pi}\int_{\RR}\sum_{j}m_j\,\frac{2\delta_j}{1+\delta_j}\,\partial_y U_j(x)\,\dd x,
\]
and using the identity
\(
\int_{\RR}\frac{2a}{(x-\gamma)^2+a^2}\,\dd x=2\pi
 \),
we obtain the closed form
\[
\boxed{\;
\mathcal{A}_{\mathrm{def}}
=\sum_{j}m_j\,\frac{\delta_j}{1+\delta_j}
=S_{\mathrm{def}}.\;}
\]
Moreover, since
\(
\frac{2\delta_j}{1+\delta_j}\partial_y U_j(x)=\frac{4\delta_j}{(x-\gamma_j)^2+(1+\delta_j)^2}
 \), it follows that \(\mathcal{A}_{\mathrm{def}}=\frac{1}{4\pi}\int_{\RR}H(x)\,\dd x\) is fully consistent with the \(c=4\) locked by the above residue unification.
This formulation provides a strict prototype of the ``horizon area law'': defect entropy is the normalized integral of the Green potential's boundary Neumann flux; the normalization factor is mathematically uniquely locked by the total boundary flux \(2\pi\) and the residue unification.
\end{proposition}

\begin{proposition}[Monotone Contraction of Information Distinguishability: Poisson Coarse-Graining Strictly Contracts Any \texorpdfstring{$f$}{f}-Divergence]
Let two defect configurations induce two nonnegative profiles \(H,\widetilde H\in L^{1}(\RR)\), normalized to probability densities
\(
h:=H/\|H\|_{1}
 \) and \(
\tilde h:=\widetilde H/\|\widetilde H\|_{1}
 \).
For any \(t>0\), let the Poisson kernel be \(P_t(x):=\frac{1}{\pi}\frac{t}{x^{2}+t^{2}}\), and define the coarse-graining
\(
h_t:=P_t*h
 \) and \(
\tilde h_t:=P_t*\tilde h
 \).
Then for any convex function \(f\) (\(f(1)=0\)), the corresponding \(f\)-divergence
\[
D_f(h\mid \tilde h):=\int_{\RR}\tilde h(x)\,f\!\left(\frac{h(x)}{\tilde h(x)}\right)\dd x
\]
satisfies the monotone contraction law
\[
\boxed{\;
D_f(h_t\mid \tilde h_t)\le D_f(h\mid \tilde h)\qquad(\forall t>0).\;}
\]
This conclusion rigorously formulates ``increasing scale leads to conceptual indistinguishability (reclamation)'' as a functional contraction law under the Markov kernel (Poisson semigroup); in particular, when \(f\) is strictly convex, unless \(h=\tilde h\) (a.e.), the inequality is strict for all \(t>0\).
\end{proposition}

\begin{proposition}[Decidable Version of the Black Hole Entropy Analogy: The Unique Normalization Making ``Deep-Defect Saturation Limit'' Strictly Equal to the Topological Count]
Let the ``area-type'' functional be
\(
A(H):=\int_{\RR}H(x)\,\dd x
 \), and consider the constant normalization
\(
S_{\lambda}:=A(H)/\lambda
 \) (\(\lambda>0\)).
If we require that for any single-atom defect \((\gamma,\delta,m)\):
\begin{enumerate}
  \item \(S_{\lambda}\) is coordinate-scale independent (invariant under Cayley family rescaling and the corresponding change of variables);
  \item \(S_{\lambda}\) is additive for mutually uncoupled defects;
  \item under the \(c=4\) normalization locked by the above residue unification, the deep-defect limit \(\delta\to\infty\) strictly saturates at the topological count \(m\);
\end{enumerate}
then the unique possible normalization is
\[
\boxed{\;\lambda=4\pi,\qquad
S_{\lambda}=\sum_{j}m_j\,p(\delta_j),\quad p(\delta)=\frac{\delta}{1+\delta}.\;}
\]
Therefore, the appearance of ``dividing by \(4\pi\)'' does not arise from an additional physical assumption, but from a purely mathematical uniqueness: the requirement that ``the normalized integral of boundary flux strictly aligns with the winding number/count in the deep-defect limit'' forces this constant and weight to be fixed.
\end{proposition}

\begin{proposition}[Ultimate Explanation of the Coefficient \(4\): Triple Unification---Curvature Scale, Residue Scale, Saturation Scale]
Within the regime framework of this paper, the coefficient \(4\) simultaneously satisfies three mutually independent and mutually derivable unification conditions:
\begin{enumerate}
  \item \textbf{Geometric side (curvature scale)}: the disk metric normalization that fixes hyperbolic curvature to \(-1\) (Proposition \ref{con:xi-curvature-normalization-alpha4});
  \item \textbf{Analytic side (residue scale)}: making the residue ratio of \(M_{H_c}\) and \(M_P\) precisely recover the dimensionless depth weight \(p(\delta)=\delta/(1+\delta)\), thereby uniquely determining \(c=4\) (given by the above residue unification);
  \item \textbf{Counting side (saturation scale)}: making the ``area-normalized entropy'' precisely saturate at the topological count in the deep-defect limit \(\delta\to\infty\) (given by the unique normalization from the above saturation requirement).
\end{enumerate}
Together, these three demonstrate that this \(4\) is a structural constant of the regime closure, not an incidental coefficient of an arithmetic object; its formal similarity with the black hole entropy formula strictly originates from the same class of ``boundary geometric quantity \(\to\) dimensionless counting quantity'' unique normalization mechanism, without introducing new dynamical interpretations.
\end{proposition}

\paragraph{Hyperbolic Area Law, Length Deficit, and Scale Fingerprint: Geometric--Analytic Alignment of Defect Entropy}

\begin{proposition}[Hyperbolic Area Law and \texorpdfstring{$4\pi$}{4pi} Normalization of Defect Entropy]
In the upper half-plane model \(\mathbb{H}\) with hyperbolic metric
\(
\dd s^{2}=(\dd x^{2}+\dd y^{2})/y^{2}
\)
(curvature \(-1\)), the area element is \(\dd A=\dd x\,\dd y/y^{2}\).
For each point defect parameter \((\gamma_j,\delta_j,m_j)\) (\(\delta_j>0\)), define the unit-width vertical strip domain
\[
R_j:=[\gamma_j-\tfrac12,\gamma_j+\tfrac12]\times[1,1+\delta_j]\subset\mathbb{H}.
\]
Then its hyperbolic area satisfies the strict identity
\[
A_{\mathbb{H}}(R_j)
=\int_{\gamma_j-1/2}^{\gamma_j+1/2}\int_{1}^{1+\delta_j}\frac{\dd y\,\dd x}{y^{2}}
=\frac{\delta_j}{1+\delta_j}.
\]
Therefore the defect entropy
\(
S_{\mathrm{def}}:=\sum_{j}m_j\,\frac{\delta_j}{1+\delta_j}
 \)
can be written as
\[
\boxed{\;
S_{\mathrm{def}}=\sum_{j}m_j\,A_{\mathbb{H}}(R_j)\;}
\]
and thus becomes the sum of hyperbolic areas of unit-width vertical strip domains.
On the other hand, for the comoving boundary profile kernel
\[
H(x)=\sum_{j} m_j\,\frac{4\delta_j}{(x-\gamma_j)^2+(1+\delta_j)^2}
\]
integrating per atom yields
\[
\int_{\RR} m_j\,\frac{4\delta_j}{(x-\gamma_j)^2+(1+\delta_j)^2}\,\dd x
=4\pi\,m_j\,\frac{\delta_j}{1+\delta_j}
=4\pi\,m_j\,A_{\mathbb{H}}(R_j),
\]
whence
\[
\boxed{\;
\frac{1}{4\pi}\int_{\RR}H(x)\,\dd x=S_{\mathrm{def}}.\;}
\]
Therefore, the essential role of the coefficient \(4\) (and the \(4\pi\) normalization) here is to strictly align the hyperbolic area (dimensionless) with the Euclidean integral of the boundary Poisson trace under the same regime normalization; its ``area law'' structure comes from the geometric--analytic unification, not from any new invariant.
\end{proposition}

\begin{proposition}[Boundary Length Deficit Formula: Entropy Is the Net Dissipation of Horocycle Length]
Under the same hyperbolic metric, the hyperbolic length of a horizontal line segment of Euclidean length \(1\) at height \(y\) is
\[
L(y):=\int_{-1/2}^{1/2}\frac{1}{y}\,\dd x=\frac{1}{y}.
\]
Thus for the above strip domain \(R_j\) we have
\[
A_{\mathbb{H}}(R_j)=\int_{1}^{1+\delta_j}\frac{1}{y^{2}}\,\dd y
=L(1)-L(1+\delta_j)
=1-\frac{1}{1+\delta_j}
=\frac{\delta_j}{1+\delta_j}.
\]
Therefore
\[
\boxed{\;
S_{\mathrm{def}}
=\sum_{j}m_j\bigl(L(1)-L(1+\delta_j)\bigr)
=\kappa-\sum_{j}\frac{m_j}{1+\delta_j}\;}
\]
where \(\kappa:=\sum_j m_j\) is the point defect index (i.e., the integer term of Proposition \ref{con:xi-index-entropy-strict-separation}).
This identity rigorously interprets the continuous entropy potential as: at the fixed base horocycle (\(y=1\)), each defect lifts the unit-width strip domain to height \(1+\delta_j\), producing a hyperbolic length deficit; entropy is the additive total of deficits.
\end{proposition}

\begin{proposition}[Optimal Feasible Range and Harmonic Mean Lower Bound of Entropy--Flux--Index]
Let the linear depth aggregate (linear flux)
\[
\Phi_{\mathrm{lin}}:=\sum_{j}m_j\,\delta_j,
\qquad
\kappa:=\sum_{j}m_j,
\]
then the defect entropy \(S_{\mathrm{def}}=\sum_j m_j\frac{\delta_j}{1+\delta_j}\) satisfies the sharp bounds
\[
\boxed{\;
\frac{\Phi_{\mathrm{lin}}}{1+\Phi_{\mathrm{lin}}}
\le
S_{\mathrm{def}}
\le
\frac{\kappa\,\Phi_{\mathrm{lin}}}{\kappa+\Phi_{\mathrm{lin}}}\;}
\]
where the upper bound is attained in the ``uniform depth'' configuration (treating \(\delta_j\) by multiplicity as \(\kappa\) numbers and setting \(\delta_j\equiv \Phi_{\mathrm{lin}}/\kappa\)), and the lower bound is the infimum for the extreme concentration configuration (one depth carrying \(\Phi_{\mathrm{lin}}\) while all others \(\downarrow 0\)).
In particular, when \(\Phi_{\mathrm{lin}}>0\), the upper bound equivalently rewrites as a harmonic mean type inequality
\[
\boxed{\;
\frac{1}{S_{\mathrm{def}}}\ge \frac{1}{\Phi_{\mathrm{lin}}}+\frac{1}{\kappa}\; }.
\]
This conclusion completely parametrizes the achievable range of the visible entropy potential within the purely point-defect class, and provides its strongest coupling constraint with the topological count \(\kappa\) and the linear flux \(\Phi_{\mathrm{lin}}\).
\end{proposition}

\begin{proposition}[Hyperbolic Exponential Law for Certificate Complexity: Finite-Order Moment Certificates Cover Only Logarithmic Radius]
Let the Euclidean boundary depth of a disk point \(a\in\mathbb{D}\) be \(h:=1-|a|\in(0,1)\), and denote its hyperbolic distance from the disk center to \(a\) (curvature \(-1\))
\[
d:=\log\frac{1+|a|}{1-|a|}.
\]
If a finite-order Toeplitz--PSD certificate (or equivalently a finite CMV certificate) can witness this point defect in the stable negative direction, then its order \(N\) must satisfy \(N\gtrsim h^{-1}\) (see Proposition \ref{con:xi-toeplitz-length-depth-duality} for the quantitative lower bound).
By the identity
\[
h=1-|a|
=\frac{2}{e^{d}+1}
\]
we obtain the equivalent hyperbolic exponential lower bound: there exists an absolute constant \(c'>0\) such that
\[
\boxed{\;
N\ \ge\ c'\,e^{d}\;}
\]
and conversely: given certificate order \(N\), the hyperbolic distance that can be stably covered is at most of order \(d\lesssim \log N\).
Therefore, the coverage radius of finite-order certificates grows only as \(\log N\) in hyperbolic geometry, while certificate resources expand exponentially with hyperbolic depth; defects closer to the boundary (larger \(d\)) necessarily enter the ``boundary layer invisible zone'' at the regime level.
\end{proposition}

\begin{proposition}[Hyperbolic Height Linearization: Entropy and ``Top Edge Length'' Form an Exponential Complement and a One-Parameter Semigroup]
Introduce the hyperbolic height variable \(s_j:=\log(1+\delta_j)\) (i.e., \(1+\delta_j=e^{s_j}\)).
Then each atom's entropy contribution satisfies
\[
\frac{\delta_j}{1+\delta_j}=1-e^{-s_j},
\]
whence
\[
S_{\mathrm{def}}=\sum_{j}m_j(1-e^{-s_j}),
\qquad
L_{\mathrm{top}}:=\sum_{j}\frac{m_j}{1+\delta_j}=\sum_{j}m_j e^{-s_j}
=\kappa-S_{\mathrm{def}}.
\]
Therefore, in hyperbolic height coordinates, entropy and top edge length form a strict exponential complementary pair: an increase in one is necessarily accompanied by an \(e^{-s}\) decay in the other.
In particular, if a uniform vertical translation \(s_j\mapsto s_j+\tau\) (\(\tau\ge0\)) is applied to all atoms, then
\[
L_{\mathrm{top}}\longmapsto e^{-\tau}L_{\mathrm{top}},
\qquad
S_{\mathrm{def}}\longmapsto \kappa-e^{-\tau}(\kappa-S_{\mathrm{def}}),
\]
so that \(\tau\) generates a one-parameter semigroup; and yields the differential identity
\[
\boxed{\;
\dd S_{\mathrm{def}}=\sum_{j}m_j e^{-s_j}\,\dd s_j\;}
\]
as the strict ledger law of ``height displacement \(\to\) entropy increment.''
\end{proposition}

\begin{proposition}[Detectability Function: The Peak Amplitude Is a Strict Quadratic Image of the Entropy Contribution]
For the single-atom profile
\(
H_{(\gamma,\delta,m)}(x)=m\,\frac{4\delta}{(x-\gamma)^2+(1+\delta)^2}
 \),
define the normalized peak amplitude
\[
a:=\frac{H_{(\gamma,\delta,m)}(\gamma)}{m}
=\frac{4\delta}{(1+\delta)^2}\in(0,1].
\]
Let the entropy contribution be \(p:=\delta/(1+\delta)\in(0,1)\), then the strict algebraic identity holds:
\[
\boxed{\;a=4p(1-p).\;}
\]
Therefore, the peak amplitude is completely determined by the entropy contribution, reaching its maximum \(a_{\max}=1\) at \(p=\tfrac12\) (equivalently \(\delta=1\)).
Conversely, given \(a\in(0,1]\), the entropy contribution has only two possible branches:
\[
p=\frac12\left(1\pm\sqrt{1-a}\right),
\]
and this two-branch ambiguity corresponds to the inherent binary uncertainty of ``amplitude alone cannot determine whether depth lies at \(\delta<1\) or \(\delta>1\)''; disambiguation must rely on phase/pole direction analytic information (such as the above residue scale).
\end{proposition}

\begin{proposition}[Complete Monotonicity of the Coarse-Grained Entropy Curve: The Scale Parameter Provides an Auditable Stieltjes Fingerprint]
Introducing the Poisson coarse-graining at scale \(t\ge0\), replacing each atomic kernel's effective width from \(1+\delta_j\) to \(t+1+\delta_j\), and defining the scale entropy
\[
S_{\mathrm{def}}(t):=\sum_{j}m_j\,\frac{\delta_j}{t+1+\delta_j}.
\]
Then \(S_{\mathrm{def}}(0)=S_{\mathrm{def}}\), and \(S_{\mathrm{def}}(t)\downarrow 0\) as \(t\to\infty\).
More strongly, \(t\mapsto S_{\mathrm{def}}(t)\) is a completely monotone function: for all integers \(n\ge0\),
\[
\boxed{\;
(-1)^n\frac{\dd^n}{\dd t^n}S_{\mathrm{def}}(t)
=n!\sum_{j}m_j\,\frac{\delta_j}{(t+1+\delta_j)^{n+1}}\ge0\; }.
\]
Therefore, \(S_{\mathrm{def}}(t)\) constitutes an auditable Stieltjes fingerprint curve; its derivative chain provides depth spectral statistics independent of the height parameter \(\gamma_j\).
Under the finite defect caliber, \(S_{\mathrm{def}}(t)\) as a rational Stieltjes function uniquely encodes \(\{(\delta_j,m_j)\}\) (height information does not appear in this scalar scale orbit), thereby further compressing the continuous entropy potential into a single-variable scale orbit.
\end{proposition}

\paragraph{Scale Entropy Curve Inversion and Two-Parameter Renormalization: Complete Compression and Minimal Coupling of Depth Spectrum}

\begin{proposition}[Stieltjes Inversion of Scale Entropy Curve: Complete Invariant of \texorpdfstring{$\{(\delta_j,m_j)\}$}{(delta,m)}]
Let point defect depth spectrum be given by \(\{(\delta_j,m_j)\}_{j=1}^{M}\), and denote
\[
a_j:=1+\delta_j>1.
\]
Define positive measure supported on \([1,\infty)\) (\(\delta_{a_j}\) being Dirac measure at \(a_j\))
\[
\mu:=\sum_{j=1}^{M} m_j\delta_j\,\delta_{a_j}.
\]
Define scale entropy curve (equivalent representation of previously defined \(S_{\mathrm{def}}(t)\) under variable \(a=1+\delta\))
\[
S(t):=\sum_{j=1}^{M}\frac{m_j\delta_j}{t+a_j}
=\int_{[1,\infty)}\frac{1}{t+a}\,\dd\mu(a),
\qquad t\ge 0.
\]
Then \(S\) is analytic on \(\mathbb{C}\setminus(-\infty,-1]\), satisfying Stieltjes inversion formula: for any continuous compactly supported test function \(\varphi\in C_c([1,\infty))\)
\[
\int_{[1,\infty)}\varphi(a)\,\dd\mu(a)
=-\lim_{\varepsilon\downarrow0}\frac{1}{\pi}\int_{1}^{\infty}\varphi(a)\,\Im S(-a+\mathrm{i}\varepsilon)\,\dd a.
\]
Therefore, the scale entropy curve \(S(t)\) (values on \(t>0\)) uniquely determines depth measure \(\mu\), thus uniquely determines depth multiset \(\{(\delta_j,m_j)\}\): recover \(a_j\) from atomic support of \(\mu\) (thus \(\delta_j=a_j-1\)), recover \(m_j\delta_j\) from atomic mass (thus \(m_j=(m_j\delta_j)/\delta_j\)).
In this sense, \(S\) is "maximal auditable compression preserving only depth information": all visible quantities independent of height coordinate \(\gamma\) and additive over defects can be viewed as functions of \(\mu\) (equivalently \(S\)).
 \end{proposition}

\begin{proposition}[Moment Closure of Finite Defects: Finite-Order Derivatives Determine All Depth Spectrum]
Let \(S\) be defined as above, and for \(n\ge0\) define
\[
\mu_n:=\frac{(-1)^n}{n!}S^{(n)}(0)
=\int_{[1,\infty)}a^{-(n+1)}\,\dd\mu(a)
=\sum_{j=1}^{M}m_j\delta_j\,a_j^{-(n+1)}.
\]
Under general position assumption that \(a_j\) are pairwise distinct, the Hankel moment matrix
\[
\mathcal{H}_N:=\bigl(\mu_{p+q}\bigr)_{0\le p,q\le N}
\]
has rank stabilizing to \(M\) for sufficiently large \(N\).
And there exists a unique (monic) polynomial
\[
Q(z)=z^{M}+q_{M-1}z^{M-1}+\cdots+q_0
\]
such that for all \(n\ge0\) there is linear recursion
\[
\sum_{r=0}^{M}q_r\,\mu_{n+r}=0,\qquad (q_M:=1),
\]
and its root set is precisely \(\{a_j^{-1}\}_{j=1}^{M}\) (counting multiplicity).
Corresponding weights \(m_j\delta_j\) are uniquely recovered by finite linear system, thus depth spectrum \(\{(\delta_j,m_j)\}\) can be uniquely determined by \(\bigl(S^{(n)}(0)\bigr)_{n=0}^{2M-1}\).
\end{proposition}

\begin{corollary}[Padé--Hankel Exact Inversion: \texorpdfstring{$2\kappa$}{2k} Order Moments Determine \texorpdfstring{$\kappa$}{k} Atomic Depth Spectrum Uniquely in Rational Class]\label{cor:xi-pade-hankel-exact-inversion}
In the finite-atom setting of the above proposition, function \(S(t)=\sum_{j=1}^{\kappa}\frac{m_j\delta_j}{t+a_j}\)'s Taylor coefficients at \(t=0\) (equivalent to moment chain \(\mu_0,\dots,\mu_{2\kappa-1}\)) uniquely determine a \((\kappa-1,\kappa)\) type Padé approximant
\[
\frac{P_{\kappa-1}(t)}{Q_{\kappa}(t)},
\qquad
\deg P_{\kappa-1}\le \kappa-1,\ \deg Q_{\kappa}\le \kappa,\ Q_{\kappa}(0)=1,
\]
matching the power series at \(t=0\) to order \(t^{2\kappa-1}\).
Then in finite defect caliber there is strict identity
\[
\boxed{\;
S(t)\ \equiv\ \frac{P_{\kappa-1}(t)}{Q_{\kappa}(t)}.\;}
\]
Thus the "pole--residue structure" of depth spectrum \(\{(a_j,m_j\delta_j)\}\) is algebraically completely determined by first \(2\kappa\) moments, with no second distinct object sharing moments in the rational Stieltjes class.
\end{corollary}
\begin{proof}
In finite-atom case \(S(t)\) is rational, denominator \(\prod_{j=1}^{\kappa}(t+a_j)\) (degree \(\kappa\)), numerator degree at most \(\kappa-1\);
and since \(a_j>1\), it is analytic in neighborhood of \(t=0\).
Therefore \(S\) itself belongs to \((\kappa-1,\kappa)\) type rational class.
Padé approximant \(\frac{P_{\kappa-1}}{Q_{\kappa}}\) is uniquely determined by "matching \(2\kappa\) Taylor coefficients";
substituting \(S\) into this matching condition shows \(S\) and its Padé approximant have identical finite-order jet at \(t=0\),
while both belong to \((\kappa-1,\kappa)\) type rational class, thus must be identically equal. \qed
\end{proof}

\begin{proposition}[``Unique Continuous Coordinatization'' of Depth Information: Linear Expression on \texorpdfstring{$S(t)$}{S(t)} Base for Visible Scalars]
Let \(\mathcal{A}\) be the class of all scalar observation functors \(F\) satisfying the following three properties:
\begin{enumerate}
  \item \textbf{Point-additive}: \(F(\mu_1+\mu_2)=F(\mu_1)+F(\mu_2)\);
  \item \textbf{Depends only on depth}: Invariant under height translation \(\gamma\mapsto\gamma+x_0\);
  \item \textbf{Covariant with additive coarse-graining semigroup}: For \(A_t:a\mapsto a+t\) (\(t\ge0\)), there exists adjoint action \(F_t\) such that \(F_t(\mu)=F((A_t)_\#\mu)\).
\end{enumerate}
Under common continuity assumption (e.g., \(F\) continuous for finite positive measures in weak topology), for any \(F\in\mathcal{A}\) there exists unique finite signed measure \(\nu_F\) such that
\[
\boxed{\;
F(\mu)=\int_{(0,\infty)} S(t)\,\dd\nu_F(t)\;}
\]
where \(S(t)=\int \frac{1}{t+a}\,\dd\mu(a)\) is the scale entropy curve defined above.
Thus in this regime class, \(S\) constitutes a "coordinatization base": any visible scalar produced only by depth and compatible with coarse-graining is a linear functional of \(S\).
\end{proposition}

\begin{proposition}[Semi-Direct Product Structure of Additive Coarse-Graining and Hyperbolic Dilation: Unique Commutation Law of Two-Parameter Renormalization Semigroup]
On depth variable \(a=1+\delta\) define two natural transformations:
\begin{enumerate}
  \item Additive coarse-graining \(A_t:a\mapsto a+t\) (\(t\ge0\));
  \item Hyperbolic dilation \(D_\tau:a\mapsto e^{\tau}a\) (\(\tau\ge0\)).
\end{enumerate}
Then they satisfy strict commutation law
\[
\boxed{\;
D_\tau\circ A_t=A_{e^{\tau}t}\circ D_\tau\;}
\]
thus generating a semi-direct product semigroup \((\tau,t)\) with composition law
\[
(\tau,t)\cdot(\tau',t')=(\tau+\tau',\,t+e^{\tau}t').
\]
This structure characterizes at regime level the "unique sustainable continuous scale freedom": all visible coarse-graining can be decomposed into unique combination of "first dilate then translate"; any additional continuous freedom incompatible with this semi-direct product degenerates to unauditable (recycled) under slice protocol.
\end{proposition}

\begin{proposition}[Completeness of Two-Parameter Analytic Generating Function: Bi-Variable Unique Determination of Depth Spectrum]
For any depth measure \(\mu\) define two-parameter generating function
\[
\mathfrak{S}(\tau,t):=\int_{[1,\infty)}\frac{1}{t+e^\tau a}\,\dd\mu(a),
\qquad (\tau,t)\in[0,\infty)^2.
\]
Then \(\mathfrak{S}\) is separately analytic in \((\tau,t)\), with scale covariance identity
\[
\boxed{\;
\mathfrak{S}(\tau,t)=e^{-\tau}\,\mathfrak{S}\!\left(0,e^{-\tau}t\right)\;}
\]
(obtained by factoring \(e^\tau\) from denominator).
For fixed \(\tau\), \(t\mapsto\mathfrak{S}(\tau,t)\) is the Stieltjes transform of pushforward measure \((D_\tau)_\#\mu\); hence \(\mathfrak{S}\)'s values on any open set uniquely determine \(\mu\) (thus determining general measure form of \(\{(\delta_j,m_j)\}\)).
\end{proposition}

\begin{proposition}[Two-Parameter Renormalization PDE: Strict Closure of Semi-Direct Action on Stieltjes Orbit]\label{prop:xi-stieltjes-semidir-pde}
Using the two-parameter generating function \(\mathfrak{S}(\tau,t)\) from above, \(\mathfrak{S}\) satisfies first-order linear PDE
\[
\boxed{\;
\partial_\tau \mathfrak{S}(\tau,t)+t\,\partial_t \mathfrak{S}(\tau,t)+\mathfrak{S}(\tau,t)=0\; }.
\]
Its characteristics are \(t e^{-\tau}=\mathrm{const}\), so \(\mathfrak{S}\) in regime closure is completely determined by unique semi-direct flow: any continuous renormalization preserving Stieltjes structure and positive weight, if incompatible with this PDE, must break complete monotonicity/additivity and be recycled in closure.
\end{proposition}
\begin{proof}
Differentiating identity \(\mathfrak{S}(\tau,t)=e^{-\tau}\mathfrak{S}(0,e^{-\tau}t)\) with respect to \(\tau\) and using chain rule yields
\(
\partial_\tau \mathfrak{S}=-\mathfrak{S}-t\,\partial_t \mathfrak{S}
\), i.e., the stated PDE. \qed
\end{proof}

\begin{proposition}[Condition Number Invariant of Depth Spectrum: Hankel Discriminant Gives Unique ``Unrecoverable Ill-Posedness Degree'']
In finite defect caliber, take \(\mu_n\) as defined above, and denote
\[
\Delta_{M-1}:=\det\bigl(\mu_{p+q}\bigr)_{0\le p,q\le M-1}.
\]
Under general position assumption that \(a_j\) are pairwise distinct, \(\Delta_{M-1}>0\), with Vandermonde-type factorization
\[
\boxed{\;
\Delta_{M-1}
=\Bigl(\prod_{j=1}^{M}m_j\delta_j\,a_j^{-1}\Bigr)\cdot
\Bigl(\prod_{1\le j<k\le M}(a_j^{-1}-a_k^{-1})^2\Bigr).\;}
\]
Therefore \(\Delta_{M-1}\) provides precise discriminant for depth spectrum "ill-posedness degree": iff there exist \(j\neq k\) with \(a_j\to a_k\) (depth collision) or \(a_j\to\infty\) (depth escape, \(a_j^{-1}\to 0\)), then \(\Delta_{M-1}\to0\).
This discriminant provides at regime level an unrecoverable ill-conditioning characterization: completely determined by depth spectrum, independent of height coordinate \(\gamma\).
\end{proposition}

\begin{proposition}[Polynomial Ill-Posedness Induced by Coarse-Graining: Additive Translation Compresses Pole Separation and Triggers Condition Number Polynomial Explosion]\label{prop:xi-stieltjes-coarsegraining-illposedness-polynomial}
In finite defect caliber, let \(S(t)=\sum_{j=1}^{\kappa}\frac{m_j\delta_j}{t+a_j}\) (\(a_j=1+\delta_j\)) and for each coarse-graining scale \(t\ge 0\) define shifted generating function
\[
S_t(u):=S(t+u)=\sum_{j=1}^{\kappa}\frac{m_j\delta_j}{u+(t+a_j)}.
\]
Its moment chain at \(u=0\) is
\[
\mu_n(t):=\frac{(-1)^n}{n!}S_t^{(n)}(0)=\sum_{j=1}^{\kappa}m_j\delta_j\,(t+a_j)^{-(n+1)}.
\]
Denote minimal separation of "inverse-width nodes"
\[
\mathrm{sep}_{\mathrm{St}}(t)
:=\min_{j\ne k}\Bigl|(t+a_j)^{-1}-(t+a_k)^{-1}\Bigr|
=\min_{j\ne k}\frac{|a_j-a_k|}{(t+a_j)(t+a_k)}.
\]
Then as \(t\to\infty\) there is deterministic decay
\[
\mathrm{sep}_{\mathrm{St}}(t)=O(t^{-2}).
\]
And for any local inversion based on finite moments \((\mu_0(t),\dots,\mu_{2\kappa-1}(t))\) (e.g., Hankel null-space method/Padé/Prony equivalent variants), its inverse map's local Lipschitz condition number \(L_{\mathrm{St}}(t)\) in general position neighborhood must satisfy lower bound
\[
\boxed{\;
L_{\mathrm{St}}(t)\ \ge\ C_\kappa\,\mathrm{sep}_{\mathrm{St}}(t)^{-(\kappa-1)}\ \gtrsim\ t^{2(\kappa-1)}\;}
\]
where constant \(C_\kappa>0\) depends only on \(\kappa\) (and chosen local norm convention), independent of height coordinates \(\{\gamma_j\}\).
Therefore, mere additive coarse-graining (translation \(a_j\mapsto t+a_j\)) already pushes depth spectrum inversion into unavoidable polynomial ill-posedness zone: scale growth compresses distinguishability and forces error amplification factor to grow at least at rate \(t^{2(\kappa-1)}\).
\end{proposition}
\begin{proof}
By algebraic identity
\(
(t+a_j)^{-1}-(t+a_k)^{-1}=(a_k-a_j)/((t+a_j)(t+a_k))
\)
we immediately get \(\mathrm{sep}_{\mathrm{St}}(t)=O(t^{-2})\).
\par
On the other hand, moment chain \(\{\mu_n(t)\}_{n=0}^{2\kappa-1}\) in general position is equivalent to a Vandermonde-type linear system with nodes
\(
z_j(t):=(t+a_j)^{-1}
\)
(the Hankel recursion/null-space construction stated in proposition "Moment Closure of Finite Defects" is its algebraic realization).
Therefore local inversion error amplification factor is unavoidably controlled by minimal node spacing: as \(\mathrm{sep}_{\mathrm{St}}(t)\downarrow 0\), any local right inverse's Lipschitz constant is at least of order \(\mathrm{sep}_{\mathrm{St}}(t)^{-(\kappa-1)}\) (power from Vandermonde determinant factor).
Substituting \(\mathrm{sep}_{\mathrm{St}}(t)=O(t^{-2})\) yields \(L_{\mathrm{St}}(t)\gtrsim t^{2(\kappa-1)}\). \qed
\end{proof}

\begin{proposition}[Exponential Law of Coarse-Graining Ill-Posedness: Scale Growth Forces Inversion Condition Number Exponential Explosion]\label{prop:xi-coarsegraining-illposedness-exponential-law}
Take sampling step size \(\Delta>0\) and let \(t\ge 0\) be additive coarse-graining scale.
For each defect \((\delta_j,\gamma_j)\) define disk point
\[
r_j(t):=\exp\!\bigl(-(t+1+\delta_j)\Delta\bigr)\,e^{-\mathrm{i}\gamma_j\Delta}\in\mathbb{D}.
\]
Denote minimal spectral spacing
\[
\mathrm{sep}(t):=\min_{j\ne k}|r_j(t)-r_k(t)|.
\]
Then there is strict scaling law
\[
\boxed{\;\mathrm{sep}(t)=e^{-t\Delta}\,\mathrm{sep}(0).\;}
\]
And for any finite-sample based local inversion algorithm (e.g., Prony/matrix pencil) in general position, Lipschitz condition number \(L(t)\) must satisfy lower bound
\[
\boxed{\;
L(t)\ \ge\ C_\kappa\,\mathrm{sep}(t)^{-(\kappa-1)}
\ \ge\ C_\kappa\,e^{(\kappa-1)t\Delta}\,\mathrm{sep}(0)^{-(\kappa-1)}\;}
\]
where \(C_\kappa>0\) depends only on \(\kappa\).
Therefore, Poisson/Lorentzian coarse-graining not only contracts distinguishability, but pushes inversion into unavoidable ill-posedness zone at strict exponential rate.
\end{proposition}
\begin{proof}
By definition \(r_j(t)=e^{-t\Delta}r_j(0)\), so \(\mathrm{sep}(t)=e^{-t\Delta}\mathrm{sep}(0)\).
On the other hand, in finite-atom caliber sample sequence has exponential sum representation \(\sum_{j=1}^{\kappa}A_j r_j(t)^n\), whose local inversion is equivalent to solving Vandermonde-type system; in general position the minimal singular value of this system is controlled by spectral spacing, so any local right inverse's Lipschitz constant is at least of order \(\mathrm{sep}(t)^{-(\kappa-1)}\) (constant depends only on \(\kappa\)).
Substituting \(\mathrm{sep}(t)\) scaling law yields exponential explosion lower bound. \qed
\end{proof}

\begin{proposition}[Scan--Certificate Equivalence Spectral Condition Number: Same Separation Synchronously Controls Hankel Ill-Posedness and Pick--Poisson Spectral Margin]\label{prop:xi-scan-toeplitz-conditioning-by-separation}
Following finite-defect setting of Proposition \ref{prop:xi-coarsegraining-illposedness-exponential-law}, for fixed \(t\ge 0\) denote disk spectral points as \(\{r_j(t)\}_{j=1}^{\kappa}\subset\mathbb{D}\).
Let \(\mathsf P(t)\in\Mat_{\kappa}(\CC)_{\mathrm{Herm}}\) be the matrix constructed from point set \(\{r_j(t)\}\) via Pick--Poisson Gram formula \eqref{eq:xi-pick-poisson-gram} (taking \(a_j=r_j(t)\) in \eqref{eq:xi-pick-poisson-gram}).
Denote endpoint weight \(p_j(t):=P_{r_j(t)}(-1)\), pseudo-hyperbolic distance \(\rho_{jk}(t):=\rho(r_j(t),r_k(t))\), and
\[
\rho_{\min}(t):=\min_{j<k}\rho_{jk}(t),\qquad
p_{\max}(t):=\max_j p_j(t),\qquad
r_{\max}(t):=\max_j|r_j(t)|.
\]
Then \(\mathsf P(t)\)'s minimal eigenvalue satisfies two types of auditable bounds:
\begin{enumerate}
  \item \textbf{(Explicit Lower Bound)} By isomorphic inequality from Corollary \ref{cor:xi-pick-poisson-gap-lower-bound},
  \[
  \boxed{\;
  \lambda_{\min}(\mathsf P(t))
  \ \ge\
  \frac{\Bigl(\prod_{j=1}^{\kappa}p_j(t)\Bigr)\Bigl(\prod_{1\le j<k\le\kappa}\rho_{jk}(t)^2\Bigr)}
  {\Bigl(\sum_{j=1}^{\kappa}p_j(t)\Bigr)^{\kappa-1}}\;}
  \]
  thus using \(\rho(a,b)=\frac{|a-b|}{|1-\overline b a|}\ge \frac{|a-b|}{1+|a||b|}\ge \frac{|a-b|}{2}\) we get
  \[
  \lambda_{\min}(\mathsf P(t))
  \ \gtrsim\
  \Bigl(\frac{\mathrm{sep}(t)}{2}\Bigr)^{\kappa(\kappa-1)}
  \cdot
  \frac{\prod_j p_j(t)}{\bigl(\sum_j p_j(t)\bigr)^{\kappa-1}}.
  \]
  \item \textbf{(Quadratic Upper Bound)} Let \(\rho_{\min}(t)=\rho(r_{j^\star}(t),r_{k^\star}(t))\).
  By \(2\times 2\) principal minor minimal eigenvalue upper bound and \(\det\mathsf P_{[j^\star,k^\star]}=p_{j^\star}(t)p_{k^\star}(t)\rho_{\min}(t)^2\) (Corollary \ref{cor:xi-pick-poisson-gap-lower-bound} \(\kappa=2\) special case) we get
  \[
  \boxed{\;
  \lambda_{\min}(\mathsf P(t))
  \ \le\
  p_{\max}(t)\,\rho_{\min}(t)^2
  \ \le\
  p_{\max}(t)\Bigl(\frac{\mathrm{sep}(t)}{1-r_{\max}(t)^2}\Bigr)^{\!2}\; }.
  \]
\end{enumerate}
Therefore when \(\mathrm{sep}(t)\downarrow 0\) (scan spectral point collision), Hankel/Vandermonde inversion condition number explodes exponentially per Proposition \ref{prop:xi-coarsegraining-illposedness-exponential-law}, while Toeplitz-side Pick--Poisson spectral margin \(\lambda_{\min}(\mathsf P(t))\) synchronously collapses (at least quadratic degradation); thus "tomography channel" and "moment certificate channel" develop consistent ill-posedness on same separation parameter.
\end{proposition}
\begin{proof}[Proof outline]
Lower bound is pointwise application of Corollary \ref{cor:xi-pick-poisson-gap-lower-bound} using \(|1-\overline b a|\le 1+|a||b|\le 2\) to control \(\rho\).
Upper bound: full matrix minimal eigenvalue does not exceed that of any \(2\times 2\) principal submatrix; for \(2\times 2\) positive definite matrix \(\lambda_{\min}\le \det/\lambda_{\max}\le 2\det/\mathrm{tr}\), using \(\mathrm{tr}\ge p_{j^\star}(t)\) yields \(\lambda_{\min}\le p_{k^\star}(t)\rho_{\min}(t)^2\le p_{\max}(t)\rho_{\min}(t)^2\).
Finally using \(|1-\overline b a|\ge 1-|a||b|\ge 1-r_{\max}(t)^2\) yields \(\rho_{\min}(t)\le \mathrm{sep}(t)/(1-r_{\max}(t)^2)\). \qed
\end{proof}

\begin{proposition}[\texorpdfstring{$\Delta$}{Delta}-Flow Moment Ladder Commutation Law: Pushforward Exponential Map Induced Derivative--Order-Raising Coupling]\label{prop:xi-delta-flow-moment-ladder-commutation}
Let depth measure \(\mu\) be supported on \([1,\infty)\).
For sampling step size \(\Delta>0\), consider Hausdorff moment generating function induced by pushforward exponential map \(a\mapsto e^{-\Delta a}\)
\[
U_0(z,\Delta):=\int_{[1,\infty)}\frac{1}{1-ze^{-\Delta a}}\,\dd\mu(a),\qquad |z|<1.
\]
And for \(k\ge 0\) define order-raising generating function
\[
U_k(z,\Delta):=\int_{[1,\infty)}\frac{a^k}{1-ze^{-\Delta a}}\,\dd\mu(a).
\]
Then for all \(|z|<1\), \(\Delta>0\) there is strict commutation law
\[
\boxed{\;
\partial_\Delta U_k(z,\Delta)\ =\ -\,z\,\partial_z U_{k+1}(z,\Delta).\;}
\]
This identity shows: evolution in step-size direction (\(\Delta\)) must inject higher depth moment \(a^{k+1}\) in the form \(z\partial_z\) at regime visible layer, thus strictly coupling "discrete sampling scale" with "higher-order depth freedom".
\end{proposition}
\begin{proof}
For \(|z|<1\) we can differentiate pointwise under integral sign:
\[
\partial_\Delta\frac{a^k}{1-ze^{-\Delta a}}
=-\,a^k\,\frac{zae^{-\Delta a}}{(1-ze^{-\Delta a})^2},
\qquad
\partial_z\frac{a^{k+1}}{1-ze^{-\Delta a}}
=a^{k+1}\,\frac{e^{-\Delta a}}{(1-ze^{-\Delta a})^2}.
\]
Multiplying both and integrating yields \(\partial_\Delta U_k=-z\partial_z U_{k+1}\). \qed
\end{proof}

\begin{proposition}[\texorpdfstring{$\Delta$}{Delta}-Flow Strict Non-Closure: Except Degenerate States Cannot Close Within Finite Moment Hierarchy]\label{prop:xi-delta-flow-nonclosure}
Following Proposition \ref{prop:xi-delta-flow-moment-ladder-commutation}.
If there exists some finite \(K\) such that for all \(|z|<1\), \(\Delta>0\), function \(U_{K+1}(z,\Delta)\) can be determined by finite set \(\{U_k(z,\Delta)\}_{k=0}^{K}\) through algebraically closed operations (finitely many additions, subtractions, multiplications, divisions, and compositions),
then \(\mu\) must be single-atom measure (support is one point).
Equivalently: for any non-degenerate depth spectrum, \(\Delta\)-flow has strict "non-closure" in regime—step evolution continuously produces new independent higher-order moment freedoms, cannot close within fixed finite moment register.
\end{proposition}
\begin{proof}[Proof outline]
By
\(
U_k(z,\Delta)=\sum_{n\ge 0}z^n\int a^k e^{-n\Delta a}\,\dd\mu(a)
\)
we know \(U_k\)'s \(z\)-coefficients give discrete Laplace moment family of \(\mu\).
Commutation law \(\partial_\Delta U_k=-z\partial_z U_{k+1}\) shows: \(\Delta\) derivative necessarily "order-raises" producing independent moment information from \(a^{k+1}\).
If some finite \(K\) closes on full domain, then these discrete Laplace moments close within finite-dimensional algebra, forcing \(\mu\)'s Laplace transform to satisfy finite-order algebraic/differential constraints; by Laplace transform uniqueness and exponential polynomial characterization, this can only occur in single-exponential case, i.e., \(\mu\) single-atom. \qed
\end{proof}

\begin{proposition}[Depth--Height Minimal Coupling: Unique Lifting from Depth Generating Function to Comoving Tomography]\label{prop:xi-depth-height-minimal-coupling}
Let complete tomographic spectral fingerprint (including height) be
\[
u_n=\sum_{j=1}^{M}4\pi\,m_j\frac{\delta_j}{1+\delta_j}\,e^{-(1+\delta_j)|n|\Delta}\,e^{-\mathrm{i}\gamma_j n\Delta},
\qquad n\in\mathbb{Z},
\]
while depth generating function (above \(S\) or this section's \(\mathfrak{S}\)) only retains \(\{(a_j,m_j\delta_j)\}\) (thus uniquely determines decay exponent \(a_j=1+\delta_j\) and amplitude weight \(m_j\delta_j/(1+\delta_j)\)).
Then in finite defect caliber, "uniquely lifting from pure depth data to height-included data" requires minimal additional information: a phase spectral carrier (phase measure) \(\sigma\) translation-covariant in \(\gamma\), such that \(e^{-\mathrm{i}\gamma_j n\Delta}\) is given by discrete spectrum of \(\sigma\).
In other words: depth generating function provides unique amplitude spectrum (decay exponent and weight), while introduction of height information must be completed at regime level through independent phase spectral carrier; no channel exists to endogenously recover \(\gamma\) from \(S(t)\) bypassing phase carrier.
This conclusion strictly formalizes "concepts compressed to single slice": depth and height form minimal direct product coupling at auditable structure level; any scheme attempting to reverse-infer height side from depth side will necessarily trigger indistinguishability (recycling) mechanism.
\end{proposition}

\begin{proposition}[Strongest Envelope of Scale Entropy Curve: Pointwise Sharp Bound Given \texorpdfstring{$(\kappa,\Phi)$}{(kappa,Phi)}]
In the above finite depth spectrum setting, denote linear flux (mass)
\[
\Phi:=\mu([1,\infty))=\sum_{j=1}^{M}m_j\delta_j,
\]
and defect index
\[
\kappa:=\sum_{j=1}^{M}m_j.
\]
Then for all \(t\ge 0\) there is uniform pointwise envelope inequality
\[
\boxed{\;
\frac{\Phi}{t+1+\Phi}\ \le\ S(t)\ \le\ \frac{\kappa\Phi}{(t+1)\kappa+\Phi}\; }.
\]
Upper bound attains equality at "uniform depth" configuration \(\delta_j\equiv \Phi/\kappa\) (for all \(j\) with \(m_j>0\)); lower bound achieves its infimum in "extreme concentration" limit: one atom carries all flux \(\Phi\) while remaining atoms have depth \(\delta\downarrow 0\).
\end{proposition}

\begin{corollary}[``Page Threshold'' Scale: Coarse-Graining Makes Scale Entropy Cross \texorpdfstring{$\kappa/2$}{kappa/2} Unique Solution and Explicit Squeezing]\label{cor:xi-scale-entropy-page-threshold}
Let scale entropy orbit \(S(t)=\sum_{j}m_j\frac{\delta_j}{t+1+\delta_j}\), denoting \(\kappa=\sum_j m_j\).
Then equation
\[
S(t)=\frac{\kappa}{2}
\]
has at most one solution \(t_\ast\) on \(t\ge 0\).
And if \(S(0)>\kappa/2\), then this solution exists uniquely, satisfying squeezing
\[
\kappa\,\frac{\delta_{\min}}{t_\ast+1+\delta_{\min}}
\ \le\ \frac{\kappa}{2}\ \le\
\kappa\,\frac{\delta_{\max}}{t_\ast+1+\delta_{\max}},
\qquad
\delta_{\min}:=\min_j\delta_j,\ \delta_{\max}:=\max_j\delta_j,
\]
thus
\[
\boxed{\;
\max\{0,\delta_{\min}-1\}\ \le\ t_\ast\ \le\ \delta_{\max}-1.\;}
\]
This threshold gives a pure-scale "information visible/invisible" boundary: once coarse-graining exceeds \(t_\ast\), scale entropy must fall into \(S(t)\le\kappa/2\) recycling phase; its position is completely sealed by depth interval, independent of height spectrum.
\end{corollary}
\begin{proof}
Since
\(
S'(t)=-\sum_j m_j\delta_j(t+1+\delta_j)^{-2}<0
\), function \(S\) is strictly decreasing on \(t\ge0\), so solution is at most unique; if \(S(0)>\kappa/2\) then by \(S(t)\to 0\) (\(t\to\infty\)) there exists unique \(t_\ast\).
By monotonicity of \(x\mapsto x/(t+1+x)\) we get
\(
\kappa\frac{\delta_{\min}}{t+1+\delta_{\min}}\le S(t)\le \kappa\frac{\delta_{\max}}{t+1+\delta_{\max}}
\);
substituting at \(t=t_\ast\) and solving inequality yields stated squeezing. \qed
\end{proof}

\begin{proposition}[Laplace--Stieltjes Factorization: Equivalent Statement of Depth Spectral Partition Function and Scale Entropy]
In finite depth spectrum setting, define depth partition function
\[
Z(u):=\sum_{j=1}^{M} m_j\delta_j\,e^{-u a_j}
\qquad (u\ge 0).
\]
Then scale entropy curve has strict Laplace factorization: for all \(t\ge 0\),
\[
\boxed{\;
S(t)=\int_{0}^{\infty}e^{-t u}\,Z(u)\,\dd u\; }.
\]
Conversely, \(Z\) is uniquely recovered from \(S\) by inverse Laplace in distribution sense; in finite defect caliber, if two scale entropy curves coincide on any set of \(t\ge 0\) with accumulation point, then corresponding partition functions \(Z\) coincide everywhere, thus depth spectrum multisets \(\{(a_j,m_j\delta_j)\}\) are identical (counting multiplicity).
\end{proposition}

\begin{proposition}[Free Energy Convexity and ``First Law'' Structure: Variance Characterization of Depth Spectrum on Single-Parameter Family]
Let \(Z(u)\) be defined as above, introducing normalized Gibbs weight
\[
p_j(u):=\frac{m_j\delta_j\,e^{-u a_j}}{\sum_{k=1}^{M}m_k\delta_k\,e^{-u a_k}}
\qquad (u>0).
\]
Denote \(\EE_u[a]:=\sum_j p_j(u)\,a_j\), \(\Var_u(a):=\sum_j p_j(u)\,(a_j-\EE_u[a])^2\).
Then log partition function satisfies strict identity
\[
\boxed{\;
-\frac{\dd}{\dd u}\log Z(u)=\EE_u[a],\qquad
\frac{\dd^2}{\dd u^2}\log Z(u)=\Var_u(a)\ge 0\; }.
\]
Therefore \(\log Z(u)\) is convex in \(u\), with second derivative precisely equal to variance of depth spectrum under Gibbs sampling; small changes in scale parameter respond only through expectation and variance in logarithmic potential.
\end{proposition}

\begin{proposition}[Moment Expansion: Full-Order Asymptotic and Depth Moment Chain of Scale Entropy at Infinity]
Denote positive moment
\[
M_n:=\int_{[1,\infty)}a^{n}\,\dd\mu(a)=\sum_{j=1}^{M}m_j\delta_j\,a_j^{n}
\qquad (n\ge 0).
\]
Then as \(t\to\infty\), scale entropy has full-order asymptotic expansion
\[
\boxed{\;
S(t)\sim \sum_{n\ge0}(-1)^n\,\frac{M_n}{t^{n+1}}\;}
\]
and for any fixed \(N\ge 0\) there is error bound
\[
S(t)=\sum_{n=0}^{N}(-1)^n\,\frac{M_n}{t^{n+1}}+O(t^{-N-2}).
\]
In particular, \(S(t)=o(t^{-1})\) iff \(\Phi=M_0=0\) (equivalent to \(\mu=0\)).
\end{proposition}

\begin{proposition}[Terminating Padé Characterization: Finite Defect iff Scale Entropy Diagonal Padé Chain Precisely Terminates in Finite Steps]
Denote \(S(t)\)'s Laurent expansion at \(t=\infty\) as \(S(t)=\sum_{n\ge0}c_n t^{-(n+1)}\).
For \(N\ge 1\), let \([N/N]\) denote diagonal Padé approximant matching first \(2N\) coefficients at \(t=\infty\): there exists unique monic polynomial \(Q_N\) (\(\deg Q_N\le N\)) and polynomial \(P_N\) (\(\deg P_N\le N-1\)) such that
\[
S(t)-\frac{P_N(t)}{Q_N(t)}=O(t^{-2N-1})\qquad (t\to\infty).
\]
Then in finite defect caliber there exists minimal integer \(N_\star\) such that \([N_\star/N_\star]\equiv S\); and under general position assumption that \(a_j\) are pairwise distinct
\[
\boxed{\;N_\star=M,\qquad Q_N(t)\equiv \prod_{j=1}^{M}(t+a_j)\ \ (N\ge M).\;}
\]
Therefore, "finite defect" can be equivalently determined as: scale entropy diagonal Padé chain precisely terminates after finite steps; terminating step number equals depth atom count \(M\) (counting multiplicity).
\end{proposition}

\begin{proposition}[Inevitable Divergence of Inversion Ill-Posedness: Hankel Discriminant Gives Unrecoverable Condition Number Lower Bound]
Under general position assumption, denote principal Hankel block
\(
H_{M-1}:=\bigl(\mu_{p+q}\bigr)_{0\le p,q\le M-1}
\)
(positive definite), with determinant \(\Delta_{M-1}\).
Then minimal eigenvalue satisfies \(\lambda_{\min}(H_{M-1})\le \Delta_{M-1}^{1/M}\), thus
\[
\boxed{\;
\|H_{M-1}^{-1}\|\ \ge\ \Delta_{M-1}^{-1/M}\; }.
\]
Therefore, when \(\Delta_{M-1}\downarrow 0\) (depth collision or depth escape), any Prony/Hankel-type inversion chain based on moment data \((\mu_n)\) (e.g., recovering \(Q\) from linear recursion and recovering \(\{a_j\}\) from its roots) unavoidably enters ill-posedness zone in regime sense: condition number diverges at least as \(\Delta_{M-1}^{-1/M}\).
\end{proposition}

\begin{proposition}[Second-Order Flux Sharp Lower Bound: Strongest Jensen--Cauchy Constraint from \texorpdfstring{$(\Phi,\Psi)$}{(Phi,Psi)} Alone]
Following notation \(a_j:=1+\delta_j>1\) and weight \(w_j:=m_j\delta_j>0\), denote
\[
\Phi:=\sum_{j=1}^{M}w_j=\mu([1,\infty)),
\qquad
\Psi:=\sum_{j=1}^{M}w_j(a_j-1)=\int_{[1,\infty)}(a-1)\,\dd\mu(a).
\]
Then for all \(t\ge 0\) there is sharp inequality
\[
\boxed{\;
S(t)\ \ge\ \frac{\Phi^2}{t\Phi+\sum_{j}w_j a_j}
\ =\ \frac{\Phi^2}{(t+1)\Phi+\Psi}\; }.
\]
And equality holds iff \(a_j\) is constant on support (i.e., all \(\delta_j\) equal, depth spectrum degenerates to single point).
Equivalently, for \(1/S(t)\) we get strongest affine upper bound
\[
\boxed{\;
\frac{1}{S(t)}\ \le\ \frac{t+1}{\Phi}+\frac{\Psi}{\Phi^2}\; }.
\]
\end{proposition}

\begin{proposition}[Superlinearity of Effective Scale Function: Monotonicity and Limit Moment Identification of \texorpdfstring{$I(t):=-S(t)/S'(t)$}{I(t)}]
For \(t>0\) define
\[
I(t):=-\frac{S(t)}{S'(t)},
\qquad
S'(t)=-\sum_{j=1}^{M}\frac{w_j}{(t+a_j)^2}<0.
\]
Then \(I\in C^\infty(0,\infty)\) satisfies strict differential inequality
\[
\boxed{\;
I'(t)=\frac{S(t)S''(t)-(S'(t))^2}{(S'(t))^2}\ \ge\ 1\; }.
\]
And \(I'(t)=1\) iff depth spectrum degenerates (all \(a_j\) equal).
Moreover as \(t\to\infty\) there is asymptotic expansion
\[
\boxed{\;
I(t)=t+\frac{\sum_{j}w_j a_j}{\Phi}+O(t^{-1})
=t+1+\frac{\Psi}{\Phi}+O(t^{-1})\; }.
\]
Therefore offset \(I(t)-t\) in large-scale limit uniquely identifies second-order flux ratio \(\Psi/\Phi\).
\end{proposition}

\begin{proposition}[Differential Discriminant of Depth Discreteness: Necessary/Sufficient Characterization of Second-Order Curvature Defect]
In above notation, define "curvature defect" functional
\[
\mathcal{D}(t):=S(t)S''(t)-2\bigl(S'(t)\bigr)^2,
\qquad t>0,
\]
where
\(
S''(t)=2\sum_{j}w_j(t+a_j)^{-3}>0
\).
Then for all \(t>0\)
\[
\boxed{\;\mathcal{D}(t)\ \ge\ 0.\;}
\]
And \(\mathcal{D}(t)\equiv 0\) (at some, thus at all \(t>0\)) iff depth spectrum degenerates to single point.
In other words, \(\mathcal{D}\) provides a pure one-dimensional in-regime discriminant: independent of height coordinate and phase carrier, determined solely by depth scale orbit, precisely characterizing "existence of multi-depth structure".
\end{proposition}

\begin{proposition}[Jacobi Certificate of Depth Spectrum: Unique Minimal Tridiagonal Register Realization]
Let depth measure \(\mu=\sum_{j=1}^{M} w_j\,\delta_{a_j}\) (\(w_j>0,a_j>1\)).
Then there exists unique symmetric tridiagonal matrix \(J\in\RR^{M\times M}\) (off-diagonal elements strictly positive) and unique constant \(c=\Phi\) such that for all \(t>0\)
\[
\boxed{\;
S(t)=c\,\langle e_1,(tI+J)^{-1}e_1\rangle\; }.
\]
where \(e_1\) is first standard basis vector.
Moreover:
\begin{enumerate}
  \item \( \mathrm{Spec}(J)=\{a_j\}\) (counting multiplicity);
  \item Denoting \(v_j\) unit eigenvector of \(J\), weights satisfy \(w_j=c\,|\langle e_1,v_j\rangle|^2\), thus tridiagonal register \((J,c)\) completely recovers \(\{(a_j,w_j)\}\) (counting multiplicity).
\end{enumerate}
Therefore, finite defect depth spectrum is not only uniquely determined by Stieltjes curve \(S\), but has canonical "minimal certificate register" realization, whose dimension precisely equals depth atom count \(M\).
\end{proposition}

\begin{proposition}[Hazard Rate Riccati Inequality: Strong Decay and Degeneracy Characterization of \texorpdfstring{$\Lambda(t):=-S'(t)/S(t)$}{Lambda(t)}]
Define scale hazard rate
\[
\Lambda(t):=-\frac{\dd}{\dd t}\log S(t)=\frac{-S'(t)}{S(t)}
=\frac{\sum_{j} w_j(t+a_j)^{-2}}{\sum_{j} w_j(t+a_j)^{-1}},
\qquad t>0.
\]
Then \(\Lambda\) is strictly monotone decreasing, satisfying Riccati-type inequality
\[
\boxed{\;
\Lambda'(t)\ \le\ -\Lambda(t)^2\; }.
\]
with equality iff depth spectrum degenerates.
Equivalently, \(\Lambda\) is controlled by strongest integrable upper bound: letting
\[
I(0^+):=\lim_{t\downarrow 0}\frac{S(t)}{-S'(t)}\in(0,\infty],
\]
then for all \(t>0\)
\[
\boxed{\;
\Lambda(t)\ \le\ \frac{1}{t+I(0^+)}\; }.
\]
This inequality gives a pure-scale "unrecoverable decay law": multi-depth structure necessarily leads to faster scale hazard rate decay than single-depth, and this accelerated decay can be witnessed by strict negativity of \(\Lambda'(t)+\Lambda(t)^2\).
\end{proposition}

\paragraph{Shadow Spectrum, Terminating \texorpdfstring{$J$}{J}-Fraction, and Bernstein Duality: Algebraic/Operator Closure Forced by Stieltjes Orbit}
In the finite-atom caliber, let the point defect depth spectrum be given by \(\{(\delta_j,m_j)\}_{j=1}^{M}\), and denote
\[
a_j:=1+\delta_j>1.
\]
Define the positive measure supported on \([1,\infty)\) (\(\delta_{a_j}\) being the Dirac measure at \(a_j\))
\[
\mu:=\sum_{j=1}^{M} m_j\delta_j\,\delta_{a_j}.
\]
Define the scale entropy curve (equivalent representation of the previously defined scale entropy curve \(S_{\mathrm{def}}(t)\) under variable \(a=1+\delta\))
\[
S(t):=\sum_{j=1}^{M}\frac{m_j\delta_j}{t+a_j}
=\int_{[1,\infty)}\frac{1}{t+a}\,\dd\mu(a),
\qquad t\ge 0.
\]
Record linear flux (mass) \(\Phi:=\sum_{j=1}^{M}m_j\delta_j\).

\begin{proposition}[``Shadow Spectrum'' Interlacing Law of Depth Spectrum: Pole--Zero Strict Interlacing of \texorpdfstring{$S(t)$}{S(t)}]\label{con:xi-shadow-spectrum-interlacing}
Let
\[
Q_{\kappa}(t):=\prod_{j=1}^{\kappa}(t+a_j),
\qquad
P_{\kappa-1}(t):=\sum_{j=1}^{\kappa} w_j\!\!\prod_{k\ne j}(t+a_k),
\]
then \(S(t)=P_{\kappa-1}(t)/Q_{\kappa}(t)\).
Under condition \(w_j>0\), \(P_{\kappa-1}\) has exactly \(\kappa-1\) distinct real roots
\[
-b_1,\dots,-b_{\kappa-1}\subset(-\infty,-1),
\]
satisfying strict interlacing
\[
\boxed{\;
a_1<b_1<a_2<b_2<\cdots<b_{\kappa-1}<a_{\kappa}\; }.
\]
This interlacing law defines a set of in-regime "shadow depths" \(\{b_j\}\): they are not equal to true depths \(\{a_j\}\), but are intermediate spectra forced by the visible Stieltjes orbit \(S\)'s rational structure between adjacent depths.
\end{proposition}
\begin{proof}
For each interval \(I_j:=(-a_{j+1},-a_j)\) (\(1\le j\le \kappa-1\)), function \(S\) is smooth and strictly decreasing on \(I_j\) (since \(S'(t)=-\sum_i w_i(t+a_i)^{-2}<0\)).
And as \(t\downarrow -a_{j+1}\), \(\frac{w_{j+1}}{t+a_{j+1}}\to +\infty\), thus \(S(t)\to +\infty\); as \(t\uparrow -a_j\), \(\frac{w_{j}}{t+a_{j}}\to -\infty\), thus \(S(t)\to -\infty\).
Hence \(S\) has exactly one zero in each \(I_j\), denoted \(-b_j\), satisfying \(-a_j>-b_j>-a_{j+1}\), i.e., \(a_j<b_j<a_{j+1}\).
Since \(P_{\kappa-1}\) has degree \(\kappa-1\), these zeros already account for all its real roots, each simple. \qed
\end{proof}

\begin{corollary}[Electrostatic Balance Characterization of Shadow Depths: Zeros Uniquely Determined by ``Splitting Charge'' Equation]\label{cor:xi-shadow-depth-electrostatic-balance}
Following Conclusion \ref{con:xi-shadow-spectrum-interlacing}, denote shadow depths as \((b_1,\dots,b_{\kappa-1})\).
Then for each \(m=1,\dots,\kappa-1\), \(b_m\in(a_m,a_{m+1})\) is the unique solution of equation
\[
\boxed{\;
\sum_{j=1}^{\kappa}\frac{w_j}{a_j-b}=0\;}
\]
on the open interval \((a_m,a_{m+1})\).
Therefore shadow depths can be interpreted as "balance points": charge terms to the left and right of \(b\) produce contributions of opposite sign, while strict monotonicity guarantees unique balance point.
\end{corollary}
\begin{proof}
By definition \(S(t)=\sum_{j=1}^{\kappa}\frac{w_j}{t+a_j}\).
Taking \(t=-b\) yields \(S(-b)=\sum_j \frac{w_j}{a_j-b}\).
Conclusion \ref{con:xi-shadow-spectrum-interlacing} already knows \(S\) is strictly decreasing on \((-a_{m+1},-a_m)\) from \(+\infty\) to \(-\infty\), so there is a unique zero in this interval; its zero point \(t=-b_m\) is equivalent to the unique solution of the stated balance equation on \((a_m,a_{m+1})\). \qed
\end{proof}

\begin{proposition}[Shadow Depth Gradient with Respect to Weights: Traction Monotonicity Completely Determined by Interval Position]\label{prop:xi-shadow-depth-gradient-wrt-weights}
Following Corollary \ref{cor:xi-shadow-depth-electrostatic-balance}.
For fixed \(m\) denote \(b_m=b_m(w)\) as the shadow depth as a function of weight vector \(w=(w_1,\dots,w_\kappa)\in\RR_{>0}^\kappa\).
Then \(b_m\) is smooth on the positive cone, satisfying explicit gradient
\[
\boxed{\;
\frac{\partial b_m}{\partial w_k}
=-\frac{1}{a_k-b_m}\Bigg/\sum_{j=1}^{\kappa}\frac{w_j}{(a_j-b_m)^2}\qquad(1\le k\le \kappa).\;}
\]
Thus: if \(a_k<b_m\) (left of shadow point), then \(\partial b_m/\partial w_k>0\); if \(a_k>b_m\) (right of shadow point), then \(\partial b_m/\partial w_k<0\).
That is, increasing left-side weight pushes shadow point rightward, while increasing right-side weight pulls shadow point leftward; this sign structure does not depend on any inversion algorithm, but is enforced solely by Stieltjes positivity and interlacing structure.
\end{proposition}
\begin{proof}
Let
\(
G(b,w):=\sum_{j=1}^{\kappa}\frac{w_j}{a_j-b}
\).
By Corollary \ref{cor:xi-shadow-depth-electrostatic-balance}, \(G(b_m(w),w)=0\).
And for \(b\in(a_m,a_{m+1})\) we have
\(
\partial_b G(b,w)=\sum_j\frac{w_j}{(a_j-b)^2}>0
\),
so by implicit function theorem differentiating with respect to \(w_k\) yields
\(
\partial_{w_k}b_m=-(\partial_{w_k}G)/(\partial_bG)
\), where \(\partial_{w_k}G(b,w)=\frac{1}{a_k-b}\).
Substituting yields the stated formula; sign conclusion follows directly from denominator positivity and sign of \(a_k-b_m\). \qed
\end{proof}

\begin{remark}[Isomorphic ``Shadow Radius'' Structure on Scan Side: Cauchy Balance Point of Hausdorff Moments]\label{rem:xi-scan-shadow-radius-cauchy-balance}
Suppose scan-side real modulus spectrum is a finite atomic measure \(\nu=\sum_{j=1}^{\kappa}c_j\,\delta_{\rho_j}\) (\(0<\rho_1<\cdots<\rho_\kappa<1\), \(c_j>0\)), and define its Cauchy transform
\(
G(x):=\sum_{j=1}^{\kappa}\frac{c_j}{x-\rho_j}
\) (\(x\in\RR\setminus\{\rho_j\}\)).
Then similarly to Corollary \ref{cor:xi-shadow-depth-electrostatic-balance}: for each \(m=1,\dots,\kappa-1\), there exists unique \(\widetilde\rho_m\in(\rho_m,\rho_{m+1})\) such that \(G(\widetilde\rho_m)=0\).
These balance points are numerator zeros of the Hausdorff moment generating function \(U(z)=\sum_{n\ge0}u_n z^n=\sum_j\frac{c_j}{1-\rho_j z}\) projected in the \(x=1/z\) variable, so the scan side is also forced to carry interlacing intermediate nodes isomorphic to depth-side shadow spectrum.
\end{remark}

\begin{proposition}[Determinant Ratio Representation of Jacobi Certificate: Shadow Spectrum is Principal Minor Spectrum]\label{con:xi-shadow-spectrum-determinant-ratio}
Let \(J\in\RR^{\kappa\times\kappa}\) be the canonical tridiagonal register given in the conclusion "Jacobi Certificate of Depth Spectrum" (symmetric with strictly positive off-diagonal elements), and let \(J^{(1)}\) denote the leading principal submatrix with first row and column deleted.
Then for all \(t>0\) there is the identity
\[
\boxed{\;
S(t)=\Phi\,\langle e_1,(tI+J)^{-1}e_1\rangle
=\Phi\,\frac{\det(tI+J^{(1)})}{\det(tI+J)}\; }.
\]
Therefore \( \mathrm{Spec}(J)=\{a_j\}\), \(\mathrm{Spec}(J^{(1)})=\{b_j\}\) (counting multiplicity), and Cauchy interlacing theorem yields the strict interlacing structure of Conclusion \ref{con:xi-shadow-spectrum-interlacing}.
\end{proposition}
\begin{proof}
By adjugate matrix identity \((A^{-1})_{11}=\mathrm{adj}(A)_{11}/\det(A)\) taking \(A=tI+J\),
\[
\langle e_1,(tI+J)^{-1}e_1\rangle
=\frac{\det\bigl((tI+J)_{(1)}\bigr)}{\det(tI+J)}
=\frac{\det(tI+J^{(1)})}{\det(tI+J)}.
\]
Multiplying by \(\Phi\) yields the conclusion. Spectral identification follows from characteristic polynomials, and interlacing between principal minor spectrum and full spectrum follows from Cauchy interlacing theorem. \qed
\end{proof}

\begin{corollary}[Multiplicative Extraction of Shadow Spectrum--Determinant Ratio: Logarithmic Ratio Invariant of Depths and Shadow Depths at \texorpdfstring{$t=0$}{t=0}]\label{cor:xi-shadow-spectrum-determinant-ratio-at-zero}
Following Conclusion \ref{con:xi-shadow-spectrum-determinant-ratio}, denote shadow depths as \((b_1,\dots,b_{\kappa-1})\).
Then at \(t=0\) there is strict multiplicative extraction law
\[
\boxed{\;
S(0)=\Phi\,\frac{\det J^{(1)}}{\det J}
=\Phi\,\frac{\prod_{j=1}^{\kappa-1}b_j}{\prod_{j=1}^{\kappa}a_j}.\;}
\]
Equivalently,
\(
\log\frac{\Phi}{S(0)}=\sum_{j=1}^{\kappa}\log a_j-\sum_{j=1}^{\kappa-1}\log b_j
\)
is a logarithmic ratio invariant completely determined by the finite certificate: it encodes both true depth and shadow depth two-layer spectral information, appearing as global spectral differential potential induced by deleting the first coordinate (rank-one externalization).
\end{corollary}
\begin{proof}
By Conclusion \ref{con:xi-shadow-spectrum-determinant-ratio} evaluating at \(t=0\) yields
\(
S(0)=\Phi\,\det(J^{(1)})/\det(J)
\).
Further using \(\det(J)=\prod_{j=1}^{\kappa}a_j\), \(\det(J^{(1)})=\prod_{j=1}^{\kappa-1}b_j\) yields the multiplicative form. \qed
\end{proof}

\begin{proposition}[Multi-Layer Shadow Spectrum Tower: Principal Minor Chain Gives Strictly Interlacing Hierarchical Horizon Lattice]\label{prop:xi-shadow-spectrum-tower-principal-minors}
Following the Jacobi certificate \(J\) of Conclusion \ref{con:xi-shadow-spectrum-determinant-ratio}.
For \(0\le r\le \kappa-1\) let \(J^{(r)}\) denote the leading principal submatrix with first \(r\) rows and columns deleted (thus \(J^{(0)}=J\) and \(J^{(\kappa-1)}\) is \(1\times 1\)).
Then there exist constants \(\Phi^{(r)}>0\) such that each layer Stieltjes orbit
\[
S^{(r)}(t):=\Phi^{(r)}\,\frac{\det(tI+J^{(r+1)})}{\det(tI+J^{(r)})}\qquad (t>0)
\]
is a rational Stieltjes function, whose pole set is \(-\mathrm{Spec}(J^{(r)})\), zero set is \(-\mathrm{Spec}(J^{(r+1)})\).
And by Cauchy interlacing theorem there is a strict hierarchical interlacing chain
\[
\boxed{\;
\mathrm{Spec}(J^{(0)}) \succ \mathrm{Spec}(J^{(1)}) \succ \cdots \succ \mathrm{Spec}(J^{(\kappa-1)}),\;}
\]
where "\(\succ\)" denotes layer-by-layer strict interlacing.
Therefore the depth spectrum not only induces one layer of shadow depths \(\{b_j\}\), but enforces an unavoidable multi-layer horizon lattice sequence; this sequence is a hierarchical invariant necessarily generated by finite register truncation and externalization projection.
\end{proposition}
\begin{proof}[Proof outline]
Applying Conclusion \ref{con:xi-shadow-spectrum-determinant-ratio} to each tail Jacobi certificate \(J^{(r)}\) (equivalent to the \(r\)-th layer tail of terminating \(J\)-fraction) yields the determinant ratio representation and pole--zero identification for each layer.
And \(J^{(r+1)}\) is a principal submatrix of \(J^{(r)}\), so its spectrum strictly interlaces with \(\mathrm{Spec}(J^{(r)})\) by Cauchy interlacing theorem. \qed
\end{proof}

\begin{proposition}[Logarithmic Potential Sum Rule: Geometric Mean Uniquely Extracted by Single-Variable Stieltjes Orbit Integral Functional]\label{prop:xi-logpotential-integral-geomean}
Following the finite-atom setting of Conclusion \ref{con:xi-shadow-spectrum-interlacing}, let total mass \(\Phi:=\sum_{j=1}^{\kappa}w_j\), and take reference orbit
\(
S_{\mathrm{ref}}(t):=\Phi/(t+1)
\).
Then the difference integral strictly converges and satisfies the identity
\[
\boxed{\;
\int_{0}^{\infty}\Bigl(S(t)-\frac{\Phi}{t+1}\Bigr)\,\dd t
=-\sum_{j=1}^{\kappa}w_j\log a_j.\;}
\]
Equivalently, letting \(p_j:=w_j/\Phi\),
\[
\boxed{\;
\int_{0}^{\infty}\Bigl(\frac{S(t)}{\Phi}-\frac{1}{t+1}\Bigr)\,\dd t
=-\sum_{j=1}^{\kappa}p_j\log a_j
=-\log\Bigl(\prod_{j=1}^{\kappa}a_j^{p_j}\Bigr).\;}
\]
Therefore the depth spectrum under regime compression still retains a "geometric mean invariant": it can be directly extracted by single-variable integral of \(S\), exhibiting standard logarithmic potential transformation law under additive coarse-graining and scale dilation.
\end{proposition}
\begin{proof}
By \(S(t)=\sum_j\frac{w_j}{t+a_j}\) we get
\[
S(t)-\frac{\Phi}{t+1}
=\sum_{j=1}^{\kappa}w_j\Bigl(\frac{1}{t+a_j}-\frac{1}{t+1}\Bigr),
\]
the right-hand side is \(O(t^{-2})\) as \(t\to\infty\), so the integral converges absolutely.
And for each \(a>1\)
\[
\int_0^\infty\Bigl(\frac{1}{t+a}-\frac{1}{t+1}\Bigr)\dd t
=\Bigl[\log(t+a)-\log(t+1)\Bigr]_{0}^{\infty}
=-\log a.
\]
Summing term by term yields the conclusion. \qed
\end{proof}

\begin{proposition}[Krein Spectral Shift Function Discretization: Staircase Shift on Shadow Spectrum Interval and Logarithmic Determinant Representation]\label{prop:xi-krein-spectral-shift-staircase-shadow}
In the setting of Conclusion \ref{con:xi-shadow-spectrum-determinant-ratio}, define normalized determinant ratio
\[
R(t):=(t+1)\,\frac{\det(tI+J^{(1)})}{\det(tI+J)}
=\frac{t+1}{\Phi}\,S(t)\qquad (t>0).
\]
Then \(R(t)>0\) and \(R(t)\to 1\) (\(t\to\infty\)).
Let staircase spectral shift function \(\xi:(1,\infty)\to\{0,1\}\) be defined by
\[
\xi(s):=\ind_{(1,a_1)}(s)\ +\ \sum_{j=1}^{\kappa-1}\ind_{(b_j,a_{j+1})}(s).
\]
Then for all \(t>0\) there is Stieltjes-type logarithmic representation
\[
\boxed{\;
\log R(t)\ =\ -\int_{(1,\infty)}\frac{\xi(s)}{t+s}\,\dd s\; }.
\]
where \(\xi\) is constant on each open interval of the interlacing lattice \(\{1,a_1,b_1,a_2,\dots,b_{\kappa-1},a_\kappa\}\), with jumps occurring only at \(\{a_j,b_j\}\).
\end{proposition}
\begin{proof}
By spectral decomposition
\(
\det(tI+J)=\prod_{j=1}^{\kappa}(t+a_j)
\),
\(
\det(tI+J^{(1)})=\prod_{j=1}^{\kappa-1}(t+b_j)
\)
we get
\[
R(t)=\frac{(t+1)\prod_{j=1}^{\kappa-1}(t+b_j)}{\prod_{j=1}^{\kappa}(t+a_j)}.
\]
On the other hand, by
\(
\int_{u}^{v}\frac{\dd s}{t+s}=\log\frac{t+v}{t+u}
\)
and definition of \(\xi\) we get
\[
\int_{(1,\infty)}\frac{\xi(s)}{t+s}\,\dd s
=\log\frac{t+a_1}{t+1}
\ +\ \sum_{j=1}^{\kappa-1}\log\frac{t+a_{j+1}}{t+b_j}
=\log\frac{\prod_{j=1}^{\kappa}(t+a_j)}{(t+1)\prod_{j=1}^{\kappa-1}(t+b_j)}.
\]
Taking the negative yields the stated identity. \qed
\end{proof}

\begin{proposition}[Uniqueness of Terminating \texorpdfstring{$J$}{J}-Fraction: Minimal Register Equivalent to Unique Positive-Coefficient Continued Fraction]\label{con:xi-terminating-jfraction-uniqueness}
There exists a unique set of coefficients
\[
\beta_0,\dots,\beta_{\kappa-1}\in(1,\infty),
\qquad
\alpha_0,\dots,\alpha_{\kappa-2}>0,
\]
such that for all \(t>0\) there is a terminating \texorpdfstring{$J$}{J}-fraction expansion
\[
\boxed{\;
S(t)
=\cfrac{\Phi}{t+\beta_0-\cfrac{\alpha_0^2}{t+\beta_1-\cfrac{\alpha_1^2}{\ddots-\cfrac{\alpha_{\kappa-2}^2}{t+\beta_{\kappa-1}}}}}\; }.
\]
And this \((\alpha,\beta)\) is termwise consistent with the Jacobi certificate: \(\beta_n=J_{n+1,n+1}\), \(\alpha_n=J_{n+1,n+2}\).
Furthermore, letting moment chain \(\mu_n:=(-1)^n S^{(n)}(0)/n!\), and defining
\[
\Delta_n:=\det(\mu_{i+j})_{0\le i,j\le n},
\qquad
\Delta_n^{(1)}:=\det(\mu_{i+j+1})_{0\le i,j\le n},
\]
the coefficients have closed-form extraction law
\[
\boxed{\;
\alpha_n^2=\frac{\Delta_{n-1}\Delta_{n+1}}{\Delta_n^2},
\qquad
\beta_n=\frac{\Delta_n^{(1)}}{\Delta_n}-\frac{\Delta_{n-1}^{(1)}}{\Delta_{n-1}},
\quad
(\Delta_{-1}:=1,\ \Delta_{-1}^{(1)}:=0)\; }.
\]
Therefore in the finite-defect case, depth visible orbit \(S\) is not only invertible, but has a unique minimal positive-coefficient register realization.
\end{proposition}
\begin{proof}[Proof outline]
By the tridiagonal structure of \(J\), repeatedly applying Schur complement to \((tI+J)\) yields the stated terminating continued fraction, with coefficient positivity given by strictly positive off-diagonal elements and positive-definiteness from \(t>0\).
Uniqueness follows from register uniqueness in the conclusion "Jacobi Certificate of Depth Spectrum".
Hankel determinant closed form is the standard extraction formula for Stieltjes moment problem: \(\Delta_n,\Delta_n^{(1)}\) give three-term recursion coefficients for orthogonal polynomials, thus recovering \((\alpha,\beta)\). \qed
\end{proof}

\begin{proposition}[Discrete Lévy Embedding of Bernstein Dual: \texorpdfstring{$1/S(t)$}{1/S(t)} has Atomic Lévy Measure Supported on Shadow Spectrum]\label{con:xi-bernstein-dual-levy-shadow-spectrum}
Let \(B(t):=1/S(t)\).
Then \(B\) belongs to the complete Bernstein class on \((0,\infty)\), and there exist unique parameters
\[
B(0)=\frac{1}{S(0)},
\qquad
\mathrm{drift}=\frac{1}{\Phi},
\]
and unique discrete Lévy measure
\(
\nu=\sum_{j=1}^{\kappa-1} c_j\,\delta_{b_j}
\)
(\(c_j>0\)) such that for all \(t>0\)
\[
\boxed{\;
B(t)=B(0)+\frac{t}{\Phi}
\;+\;\int_{(0,\infty)}\frac{t}{t+s}\,\dd\nu(s)
=B(0)+\frac{t}{\Phi}+\sum_{j=1}^{\kappa-1}c_j\,\frac{t}{t+b_j}\; }.
\]
where \(\{b_j\}\) are precisely the shadow depths (also \(\mathrm{Spec}(J^{(1)})\)) from Conclusion \ref{con:xi-shadow-spectrum-interlacing}, and weights satisfy explicit residue formula
\[
\boxed{\;
c_j=-\frac{Q_{\kappa}(-b_j)}{b_j\,P'_{\kappa-1}(-b_j)}\ >\ 0\; }.
\]
This conclusion shows: the dual coarse-graining residual of depth orbit is not supported on true depths \(\{a_j\}\) at external interface, but on shadow spectrum \(\{b_j\}\) forcibly generated by regime.
\end{proposition}
\begin{proof}[Proof outline]
By \(S(t)=P_{\kappa-1}(t)/Q_{\kappa}(t)\) we get \(B(t)=Q_{\kappa}(t)/P_{\kappa-1}(t)\) as rational function.
Conclusion \ref{con:xi-shadow-spectrum-interlacing} gives zeros of \(P_{\kappa-1}\) as \(\{-b_j\}\) and distinct, so \(B\) has simple poles at \(\{-b_j\}\) and admits partial fraction decomposition.
As \(t\to\infty\), \(S(t)\sim \Phi/t\), so \(B(t)\sim t/\Phi\) gives drift.
Rewriting partial fraction in kernel form \(t/(t+b_j)=1-b_j/(t+b_j)\) yields the stated Lévy representation; residue comparison gives explicit formula for \(c_j\). \qed
\end{proof}

\begin{proposition}[Two-Moment Envelope Theorem: Sharpest Pointwise Bound Given Only \texorpdfstring{$(\Phi,S(0))$}{(Phi,S0)}]\label{con:xi-two-moment-envelope}
Let \(\Phi=\sum_j w_j\) and \(S(0)=\sum_j w_j/a_j\) be given.
Then for any finite depth spectrum satisfying these two-moment constraints, for all \(t\ge 0\) there is sharp envelope
\[
\boxed{\;
\frac{S(0)}{1+t}\ \le\ S(t)\ \le\ \frac{\Phi\,S(0)}{\Phi+tS(0)}\; }.
\]
Upper bound attains equality if and only if depth spectrum degenerates to single point (all \(a_j\) equal); lower bound is achieved in closure sense by "boundary-layer--deep-layer two-phase" limit.
\end{proposition}
\begin{proof}
Since \(a_j\ge 1\), we have \(t+a_j\le (1+t)a_j\), thus
\(
S(t)=\sum_j w_j/(t+a_j)\ge \frac{1}{1+t}\sum_j w_j/a_j=\frac{S(0)}{1+t}
\).
On the other hand, Conclusion \ref{con:xi-bernstein-dual-levy-shadow-spectrum} gives
\(
B(t)=1/S(t)\ge B(0)+t/\Phi=1/S(0)+t/\Phi
\);
taking reciprocal yields upper bound. \qed
\end{proof}

\begin{proposition}[Scale Cauchy--Schwarz Constraint: Universal Quadratic Inequality of \texorpdfstring{$S(t)$}{S(t)} and \texorpdfstring{$-S'(t)$}{-S'(t)}]\label{con:xi-scale-cauchy-schwarz}
For any \(t\ge 0\) there holds
\[
\boxed{\;
S(t)^2\ \le\ \Phi\,\bigl(-S'(t)\bigr),
\qquad
-S'(t)=\sum_{j=1}^{\kappa}\frac{w_j}{(t+a_j)^2}\; }.
\]
Equality holds if and only if depth spectrum degenerates to single point.
\end{proposition}
\begin{proof}
By Cauchy--Schwarz inequality
\[
\Bigl(\sum_j \frac{w_j}{t+a_j}\Bigr)^2
=\Bigl(\sum_j \sqrt{w_j}\cdot\frac{\sqrt{w_j}}{t+a_j}\Bigr)^2
\le \Bigl(\sum_j w_j\Bigr)\Bigl(\sum_j \frac{w_j}{(t+a_j)^2}\Bigr)
=\Phi\,(-S'(t)).
\]
Equality iff vectors \(\sqrt{w_j}\) and \(\sqrt{w_j}/(t+a_j)\) are proportional, i.e., all \(t+a_j\) equal, thus all \(a_j\) equal. \qed
\end{proof}

\begin{proposition}[Alternating Bounds of Continued Fraction Convergents: Finite Register Gives Full-Scale Uniform Upper/Lower Certificates]\label{con:xi-jfraction-convergents-alternating-bounds}
Let convergents of the terminating \texorpdfstring{$J$}{J}-fraction from Conclusion \ref{con:xi-terminating-jfraction-uniqueness} be \(S_n(t)\) (by convention "truncated to \(n\)-th layer"), written in reduced form \(S_n(t)=P_n(t)/Q_n(t)\).
Then for all \(t>0\) there is strict alternating squeezing chain
\[
\boxed{\;
S_0(t)\ \ge\ S_2(t)\ \ge\ \cdots\ \ge\ S(t)\ \ge\ \cdots\ \ge\ S_3(t)\ \ge\ S_1(t)\; }.
\]
And for \(0\le n\le \kappa-2\) there is continuant determinant identity
\[
\boxed{\;
S_{n+1}(t)-S_n(t)=
\frac{(-1)^n\,\Phi\,\prod_{k=0}^{n}\alpha_k^2}{Q_{n+1}(t)\,Q_{n}(t)}\; }.
\]
Therefore, finite \((\alpha_k,\beta_k)\) alone give uniform upper/lower bounds and explicit error width across all scales \(t>0\).
\end{proposition}

\begin{remark}[``Optimal Compression Node'' Meaning of Shadow Spectrum: Extremal Envelope, Atomic Lévy Measure, and Interlacing Error Lattice]\label{rem:xi-shadow-spectrum-optimal-compression-nodes}
Conclusion \ref{con:xi-jfraction-convergents-alternating-bounds} shows: under regime allowing access only to finite register coefficients \((\alpha_k,\beta_k)\), continued fraction convergents give full-scale strongest upper/lower certificates for true orbit \(S(t)\); error width is explicitly given by continuant identity, contracting deterministically with register length.
On the other hand, discrete Lévy measure of Bernstein dual \(B(t)=1/S(t)\) is forcibly pinned on shadow spectrum \(\{b_j\}\) (Conclusion \ref{con:xi-bernstein-dual-levy-shadow-spectrum}), thus "external coarse-graining residual" atomic nodes are unique in algebraic sense.
\par
Their combination gives a strict optimality interpretation: shadow spectrum is not an arbitrarily introduced intermediate parameter, but under Stieltjes positivity and interlacing structure, is the minimal node set unavoidable for maintaining order-preserving full-scale certificates (real negative poles, positive residues, alternating errors); any attempt to compress/approximate \(S\) and \(1/S\) without externalizing new continuous degrees of freedom will be forced back to the rational structure generated by interlacing lattice \(\{-a_j,-b_j\}\).
\end{remark}

\begin{proposition}[Structural Stability of Shadow Spectrum: Regular Response and Unavoidable Collision Ill-Posedness Under Weight Perturbation]\label{con:xi-shadow-spectrum-stability}
Under condition that \(a_1<\cdots<a_\kappa\) are fixed and distinct, shadow depths \((b_1,\dots,b_{\kappa-1})\) as function of normalized weight vector
\(
\tilde w:=w/\Phi\in\{x\in\RR_{>0}^\kappa:\sum x_j=1\}
\)
are analytic at interior points of the simplex.
And denoting
\(
\mathrm{sep}:=\min_{1\le j\le \kappa-1}(a_{j+1}-a_j)>0
\),
on any fixed compact interior domain there exists constant \(C=C(\{a_j\})>0\) such that
\[
\max_{1\le j\le \kappa-1}\bigl|\delta b_j\bigr|
\ \le\ C\,\bigl|\delta \tilde w\bigr|_1.
\]
When \(\mathrm{sep}\downarrow 0\) (adjacent depth collision), the above Lipschitz constant must diverge at rate \(\mathrm{sep}^{-1}\), thus shadow spectrum simultaneously provides "regular response" and "regime ill-posedness zone" necessary/sufficient signals.
\end{proposition}

\begin{proposition}[Operator Meaning of Shadow Spectrum: External Lévy Atoms are Spectral Flow Carriers of Leading Principal Submatrix]\label{con:xi-shadow-spectrum-operator-meaning}
Combining Conclusion \ref{con:xi-shadow-spectrum-determinant-ratio} and Conclusion \ref{con:xi-bernstein-dual-levy-shadow-spectrum}, the discrete Lévy measure of \(B(t)=1/S(t)\)
\(
\nu=\sum_{j=1}^{\kappa-1}c_j\,\delta_{b_j}
\)
satisfies:
\[
\mathrm{supp}(\nu)=\{b_j\}=\mathrm{Spec}(J^{(1)}),
\]
and taking unit eigenvector \(u_j\) of \(J^{(1)}\) (corresponding to eigenvalue \(b_j\)), then \(c_j\) is proportional to "first-coordinate observation" coupling strength and can be extracted by residue closed form:
\[
c_j=-\frac{Q_{\kappa}(-b_j)}{b_j\,P'_{\kappa-1}(-b_j)}.
\]
Therefore, dual coarse-graining residual of depth orbit does not directly remember true depths, but forms auditable spectral flow carrier at external interface through leading principal submatrix modes; shadow spectrum thereby obtains strict external operator interpretation.
\end{proposition}

\begin{proposition}[Support Squeezing: \texorpdfstring{$I(t)=-S(t)/S'(t)$}{I(t)} Squeezes Depth Interval at Every Scale]\label{con:xi-support-squeeze-I}
Let
\[
I(t):=-\frac{S(t)}{S'(t)}=\frac{S(t)}{-S'(t)}\qquad (t>0),
\]
and denote \(a_{\min}:=\min_j a_j\), \(a_{\max}:=\max_j a_j\).
Then for all \(t>0\)
\[
\boxed{\;
t+a_{\min}\ \le\ I(t)\ \le\ t+a_{\max}\; }.
\]
and equality holds if and only if depth spectrum degenerates to single point.
In particular, \(I(t)-t\) as a pure one-dimensional visible quantity gives auditable squeezing of true depth support interval without any phase carrier input.
\end{proposition}
\begin{proof}
Let \(x_j:=t+a_j\), and define probability weight
\(
p_j(t):=\frac{w_j x_j^{-2}}{\sum_{k}w_k x_k^{-2}}
\).
Then
\[
I(t)=\frac{\sum_j w_j x_j^{-1}}{\sum_j w_j x_j^{-2}}
=\sum_{j}p_j(t)\,x_j,
\]
being a convex combination of \(\{x_j\}\), thus lies between \(\min_j x_j=t+a_{\min}\) and \(\max_j x_j=t+a_{\max}\).
Equality holds iff all \(x_j\) equal, i.e., all \(a_j\) equal. \qed
\end{proof}

\endinput

\paragraph{Truncated moment certificates and Markov--Stieltjes extremals: uncertainty intervals under finite register length and endpoint realization}
Let \(\mu\) be a finite positive measure supported on \([1,\infty)\), and let its Stieltjes orbit be
\[
S(t):=\int_{[1,\infty)}\frac{1}{t+a}\,\dd\mu(a),\qquad t\ge 0.
\]
Since the spectral gap \(a\ge 1\), the function \(S\) is analytic in a neighborhood of \(t=0\) and induces the truncated moment chain (namely, the Taylor coefficients of \(S\) at \(t=0\))
\[
\mu_n:=\frac{(-1)^n}{n!}S^{(n)}(0)=\int_{[1,\infty)}a^{-(n+1)}\,\dd\mu(a),\qquad n\ge 0.
\]

\begin{proposition}[Markov--Stieltjes extremal principle: strongest feasible interval under truncated moment information and endpoint realization]\label{prop:xi-markov-stieltjes-extremal-principle}
Fix an integer \(N\ge 1\) and given a truncated moment vector \((\mu_0,\dots,\mu_{2N-1})\).
Let \(\mathfrak{M}_N\) be the set of all finite positive measures satisfying the above \(2N\) moment constraints:
\[
\mathfrak{M}_N:=\Bigl\{\nu\ \text{is a finite positive measure on \([1,\infty)\)}:\ \int a^{-(n+1)}\,\dd\nu(a)=\mu_n\ (0\le n\le 2N-1)\Bigr\}.
\]
Then for any fixed \(t>0\), the feasible value set
\[
\bigl\{S_\nu(t):=\int\frac{1}{t+a}\,\dd\nu(a)\ :\ \nu\in\mathfrak{M}_N\bigr\}
\]
is a closed interval \([S_N^{-}(t),S_N^{+}(t)]\).
Here the endpoint functions \(S_N^\pm\) are two Stieltjes-type rational functions uniquely determined by \((\mu_0,\dots,\mu_{2N-1})\), and can be taken as adjacent convergents of the same positive-coefficient Stieltjes continued fraction (even/odd extremal branches).
Furthermore, the extremal measures attaining the upper and lower endpoints must be atomic measures with at most \(N\) atoms.
\end{proposition}
\begin{proof}[Proof outline]
This is the standard extremal conclusion for the truncated Stieltjes moment problem (Markov theorem/endpoint case of Ne\-van\-lin\-na parametrization).
For fixed \(t>0\), the mapping \(\nu\mapsto S_\nu(t)\) is a continuous linear functional; by Nevanlinna parametrization, the set of all feasible values is precisely a closed interval, whose endpoints are realized by extreme solutions.
The extreme solutions of the truncated moment problem must be discrete measures with at most \(N\) atoms, and correspond to the two adjacent convergents (even/odd extremal branches) of the positive-coefficient continued fraction constructed by the Stieltjes algorithm from \((\mu_0,\dots,\mu_{2N-1})\). This completes the proof.
\end{proof}

\begin{proposition}[Minimum atomic realization of truncated moments: provable lower bound of defect degree and Gauss-type reconstruction]\label{prop:xi-truncated-moment-min-atomic-realization}
Under the setting of Proposition \ref{prop:xi-markov-stieltjes-extremal-principle}, denote the Hankel matrix block
\[
H_{N-1}:=\bigl(\mu_{i+j}\bigr)_{0\le i,j\le N-1}\succeq 0,
\qquad
\mathfrak{d}_N:=\mathrm{rank}(H_{N-1})\le N.
\]
Then there exist a set of nodes \(\{a_j^{(N)}\}_{j=1}^{\mathfrak{d}_N}\subset[1,\infty)\) and weights \(\{w_j^{(N)}\}_{j=1}^{\mathfrak{d}_N}\subset(0,\infty)\) such that
\[
\mu_n=\sum_{j=1}^{\mathfrak{d}_N} w_j^{(N)}\,\bigl(a_j^{(N)}\bigr)^{-(n+1)},
\qquad n=0,1,\dots,2N-1.
\]
Moreover, this minimum atomic realization can be constructed from the orthogonal polynomial system generated by the truncated moments: the nodes are the real zeros (lying in \((1,\infty)\)) of the corresponding orthogonal polynomial, and the weights are determined in closed form by the Christoffel numbers.
Therefore, the "defect degree" has a truncated version beyond the full-information limit \(\kappa\): given a finite certificate length \(2N\), the minimum realizable defect degree within the regime is at most \(N\), and can be constructively recovered from purely algebraic objects.
\end{proposition}
\begin{proof}[Proof outline]
From the truncated moments \((\mu_0,\dots,\mu_{2N-1})\) one can define a positive definite/semi-definite inner product on \(\mathrm{span}\{1,a^{-1},\dots,a^{-N}\}\), and obtain orthogonal polynomials and three-term recurrence; the spectrum of the corresponding truncated Jacobi matrix gives the nodes, while the first coordinate weights give the Christoffel numbers.
\(\mathfrak{d}_N=\mathrm{rank}(H_{N-1})\) is precisely the minimum atomic number lower bound for this finite-dimensional moment problem, and attainability is given by Gauss-type quadrature/spectral measure construction. This completes the proof.
\end{proof}

\begin{proposition}[Depth spectral Hankel determinant closed form: Vandermonde square law and algebraic closure of separation pathology]\label{prop:xi-depth-hankel-determinant-vandermonde-square}
Under the finite atomic regime, suppose
\[
\mu=\sum_{j=1}^{\kappa}w_j\,\delta_{a_j},
\qquad
1<a_1<a_2<\cdots<a_\kappa,\qquad w_j>0,
\]
and let the moment chain \(\mu_n=\int a^{-(n+1)}\,\dd\mu(a)\).
Denote \(b_j:=a_j^{-1}\in(0,1)\), \(\widetilde w_j:=w_j b_j>0\), then
\[
\mu_n=\sum_{j=1}^{\kappa}\widetilde w_j\,b_j^{n}\qquad(n\ge 0).
\]
Let the Hankel principal block \(\mathsf H_{\kappa-1}:=(\mu_{p+q})_{0\le p,q\le \kappa-1}\).
Then the Vandermonde factorization holds
\[
\boxed{\;
\mathsf H_{\kappa-1}=V\,\mathrm{diag}(\widetilde w_1,\dots,\widetilde w_\kappa)\,V^{\mathsf T},\qquad
V_{pj}:=b_j^{p}\;}
\]
and thus the determinant closed form
\begin{equation}\label{eq:xi-depth-hankel-det-vandermonde}
\boxed{\;
\det \mathsf H_{\kappa-1}
=\left(\prod_{j=1}^{\kappa}\widetilde w_j\right)
\left(\prod_{1\le j<k\le \kappa}(b_k-b_j)^2\right).\;}
\end{equation}
In particular, \(\det \mathsf H_{\kappa-1}=0\) if and only if there exists a spectral collision \(b_j=b_k\) or weight degeneration \(\widetilde w_j=0\).
Moreover, any local stability constant based on Hankel inversion (such as Prony/Pad\'e/Hankel null-space methods) cannot remain uniformly bounded at the \(b\)-collision limit: when \(\prod_{j<k}|b_k-b_j|\downarrow 0\), the error amplification factor diverges at least at a rate \(\bigl(\prod_{j<k}|b_k-b_j|^{-2}\bigr)\).
\end{proposition}

\begin{proof}
From
\(
\mu_{p+q}=\sum_j \widetilde w_j\,b_j^{p+q}
\)
we directly obtain the matrix factorization \(\mathsf H_{\kappa-1}=V\,\mathrm{diag}(\widetilde w)\,V^{\mathsf T}\).
Taking determinants and using \(\det(V)=\prod_{j<k}(b_k-b_j)\) yields \eqref{eq:xi-depth-hankel-det-vandermonde}.
The last statement is a standard condition number lower bound: \(\|\mathsf H_{\kappa-1}^{-1}\|\) is at least of the same scale as \((\det\mathsf H_{\kappa-1})^{-1}\) (up to a polynomial factor controlled by \(\|\mathsf H_{\kappa-1}\|\)), so collapse of the Vandermonde square factor must trigger inversion pathology explosion. This completes the proof.
\end{proof}

\begin{proposition}[Hankel determinant covariance law under scaling renormalization: scale factors can be explicitly extracted]\label{prop:xi-hankel-det-scaling-covariance}
Under the finite atomic setting of Proposition \ref{prop:xi-depth-hankel-determinant-vandermonde-square}, write \(a_j=1+\delta_j\) (\(\delta_j>0\)).
For any scaling parameter \(m>0\), define the scaled depth nodes
\[
\widetilde a_j:=1+m\delta_j,
\qquad
\widetilde b_j:=\widetilde a_j^{-1}=\frac{1}{1+m\delta_j},
\qquad
\widetilde{\widetilde w}_j:=w_j\widetilde b_j,
\]
and let \(\widetilde{\mathsf H}_{\kappa-1}\) be the Hankel principal block induced by the scaled moment chain
\(
\widetilde\mu_n:=\sum_{j=1}^{\kappa}w_j\,\widetilde a_j^{-(n+1)}
=\sum_{j=1}^{\kappa}\widetilde{\widetilde w}_j\,\widetilde b_j^{n}
\).
Then there exists an explicit covariance factor \(J_{\kappa-1}(m;\delta)>0\) such that
\[
\boxed{\;
\det \widetilde{\mathsf H}_{\kappa-1}
=\det \mathsf H_{\kappa-1}\cdot J_{\kappa-1}(m;\delta),\;}
\]
and one can take
\[
\boxed{\;
J_{\kappa-1}(m;\delta)
=m^{\kappa(\kappa-1)}
\prod_{j=1}^{\kappa}\Bigl(\frac{1+\delta_j}{1+m\delta_j}\Bigr)^{2\kappa-1}.\;}
\]
Therefore the scaling flow does not change the algebraic degeneration type of "collision/merging": when there exists \(\delta_j\to\delta_k\) (equivalent to \(b_j\to b_k\)), both \(\det \mathsf H_{\kappa-1}\) and \(\det\widetilde{\mathsf H}_{\kappa-1}\) exhibit second-order vanishing by the Vandermonde square law (Proposition \ref{prop:xi-depth-hankel-determinant-vandermonde-square}), while the scaling only introduces separable single-point scale factors through \(J_{\kappa-1}\).
\end{proposition}
\begin{proof}[Proof outline]
By \eqref{eq:xi-depth-hankel-det-vandermonde}, for scaled nodes we have
\[
\det \widetilde{\mathsf H}_{\kappa-1}
=\Bigl(\prod_{j=1}^{\kappa}w_j\widetilde b_j\Bigr)\Bigl(\prod_{1\le j<k\le\kappa}(\widetilde b_k-\widetilde b_j)^2\Bigr).
\]
Now
\(
\widetilde b_k-\widetilde b_j=\frac{1}{1+m\delta_k}-\frac{1}{1+m\delta_j}
=\frac{m(\delta_j-\delta_k)}{(1+m\delta_j)(1+m\delta_k)}
\).
Comparing this with \(b_k-b_j=\frac{\delta_j-\delta_k}{(1+\delta_j)(1+\delta_k)}\) and combining the products over all pairs \((j,k)\), we obtain the stated \(J_{\kappa-1}(m;\delta)\). This completes the proof.
\end{proof}

\begin{proposition}[Order preservation of truncated Pad\'e--Stieltjes approximation: poles are real negative, residues are positive, and yield monotone convergence across all scales]\label{prop:xi-pade-stieltjes-order-preserving}
If \(S\) is a Stieltjes function as above, then its diagonal Pad\'e approximant \([N/N]\) at \(t=0\) (matching the first \(2N\) Taylor coefficients) is still a Stieltjes-type rational function: all its poles lie in \((-\infty,-1)\) and are simple, and the residues in the corresponding partial fraction expansion are strictly positive.
Moreover, for any \(t>0\), the sequence \([N/N](t)\) splits into two monotone chains by parity in \(N\), squeezing the true value \(S(t)\) from above and below; and as \(N\to\infty\) the squeezing interval converges to the singleton \(\{S(t)\}\).
This conclusion elevates the auditability of "finite register certificates" to order preservation: the approximation process generates neither spurious poles nor sign-flipping spurious weights, thus excluding at the regime level any non-physical defect indices introduced by approximation.
\end{proposition}
\begin{proof}[Proof outline]
The fact that Pad\'e approximants of the Stieltjes class preserve positive weights and real negative pole structure is a direct consequence of continued fraction/orthogonal polynomial theory; the parity monotonicity corresponds to the alternating squeezing property of Stieltjes continued fraction convergents, and convergence comes from complete monotonicity and compactness control of the Stieltjes transform on \((0,\infty)\). This completes the proof.
\end{proof}

\begin{proposition}[Pure tail identifiability of the shallowest depth atom: moment ratio limit determines \texorpdfstring{$a_{\min}$}{a_min} and gives exponential convergence rate]\label{prop:xi-moment-ratio-identifies-amin}
Under the finite atomic regime, suppose
\(
\mu=\sum_{j=1}^{\kappa} w_j\,\delta_{a_j}
\)
with \(1<a_1:=a_{\min}<a_2\le\cdots\le a_{\kappa}\), \(w_j>0\).
Let the moment chain
\[
\mu_n:=\int a^{-(n+1)}\,\dd\mu(a)=\sum_{j=1}^{\kappa} w_j\,a_j^{-(n+1)}\qquad(n\ge 0).
\]
Then the limits exist
\[
\boxed{\;
\lim_{n\to\infty}\frac{\mu_n}{\mu_{n+1}}=a_1,
\qquad
\lim_{n\to\infty}\frac{\mu_{n+1}}{\mu_n}=\frac{1}{a_1}\; }.
\]
Moreover, the convergence is exponential: there exists a constant \(C>0\) such that
\[
\boxed{\;
\left|\frac{\mu_n}{\mu_{n+1}}-a_1\right|
\ \le\ C\left(\frac{a_1}{a_2}\right)^{n}\qquad(n\to\infty)\; }.
\]
Therefore the shallowest depth \(\delta_{\min}=a_1-1\), under the condition of not introducing any phase carrier and not performing full spectral inversion, can be uniquely recovered from the ratio limit of the single-variable moment chain tail; the spectral gap ratio \(a_1/a_2\) precisely controls its regime-level identifiability speed.
\end{proposition}
\begin{proof}
Write the moment chain as \(\mu_n=w_1a_1^{-(n+1)}(1+R_n)\), where
\(
R_n:=\sum_{j=2}^{\kappa}\frac{w_j}{w_1}\bigl(\frac{a_1}{a_j}\bigr)^{n+1}
\).
Then \(|R_n|\le C_0(\frac{a_1}{a_2})^{n}\) (for some constant \(C_0>0\)), and
\[
\frac{\mu_n}{\mu_{n+1}}
=a_1\cdot\frac{1+R_n}{1+R_{n+1}}
=a_1\left(1+\frac{R_n-R_{n+1}}{1+R_{n+1}}\right).
\]
Since \(R_n-R_{n+1}=O((a_1/a_2)^n)\) we obtain the stated exponential error bound. This completes the proof.
\end{proof}

\begin{proposition}[Tail rate limit law of Laplace partition: \texorpdfstring{$\delta_{\min}$}{delta_min} is the rate of the partition function]\label{prop:xi-partition-tail-rate-deltamin}
Under the finite atomic setting of the previous proposition, let \(\delta_j:=a_j-1\), and define the depth partition function
\[
Z_\delta(u):=\sum_{j=1}^{\kappa} w_j\,e^{-u\delta_j}\qquad(u\ge 0).
\]
Then the limit exists
\[
\boxed{\;
\lim_{u\to\infty}-\frac{1}{u}\log Z_\delta(u)=\delta_{\min}:=\min_j \delta_j\; }.
\]
Moreover, if \(\delta_2\) is the second smallest depth, then there is an exponential error bound: there exists a constant \(C>0\) such that
\[
\left|-\frac{1}{u}\log Z_\delta(u)-\delta_{\min}\right|
\ \le\ \frac{1}{u}\log\!\left(1+C\,e^{-u(\delta_2-\delta_{\min})}\right).
\]
This conclusion elevates the identifiability of "shallowest depth" to a thermodynamic limit law: \(\delta_{\min}\) is the rate function of the partition function, and the depth spectral gap \(\delta_2-\delta_{\min}\) determines the exponential scale of its convergence.
\end{proposition}
\begin{proof}
Write
\(
Z_\delta(u)=e^{-u\delta_{\min}}\bigl(w_1+\sum_{j\ge 2} w_j e^{-u(\delta_j-\delta_{\min})}\bigr)
\),
take logarithm and divide by \(u\) to obtain the limit.
The error bound is given by the exponential decay of the second term inside the bracket. This completes the proof.
\end{proof}

\begin{corollary}[Minimum observation scale of tail rate certificate: \texorpdfstring{$\delta_{\min}$}{delta_min}-driven indistinguishability lower bound]\label{cor:xi-partition-tail-min-scale-lower-bound}
Still under the finite defect setting of Proposition \ref{prop:xi-partition-tail-rate-deltamin}, suppose the regime only allows observation of the partition tail rate \(Z_\delta(u)\) (or equivalently observation of \(-\frac{1}{u}\log Z_\delta(u)\)) on a discrete grid
\(
u\in\{0,\Delta u,2\Delta u,\dots,M\Delta u\}
\),
and the observation accuracy requirement is to achieve relative error on the order of \(\varepsilon\in(0,1)\).
Then to distinguish in the worst case between "critical defect-free (\(\delta_{\min}=0\) in the closure sense)" and "existence of shallowest depth \(\delta_{\min}>0\)", the required maximum observation scale must satisfy
\[
\boxed{\;
M\Delta u\ \gtrsim\ \frac{1}{\delta_{\min}}\log\!\frac{1}{\varepsilon}\;}
\]
where the implicit constant depends only on a priori bounds on weight ratios (e.g., an upper bound on \(w_2/w_1\)).
This lower bound characterizes the invisible boundary layer within the regime: the smaller \(\delta_{\min}\) is, the larger \(u\) (longer "regime time") one must push the observation to in order to produce auditable separation.
\end{corollary}
\begin{proof}[Proof outline]
By \(Z_\delta(u)\ge w_1e^{-u\delta_{\min}}\) and \(Z_\delta(u)\le (w_1+\cdots+w_\kappa)e^{-u\delta_{\min}}\) we have: if \(\delta_{\min}=0\) (closure regime) then \(Z_\delta(u)\ge w_1\) does not decay with \(u\); whereas if \(\delta_{\min}>0\) then \(Z_\delta(u)\le \Phi\,e^{-u\delta_{\min}}\) decays exponentially (\(\Phi:=\sum_j w_j\)).
Therefore to distinguish the two under relative error \(\varepsilon\), one must enter the exponential tail regime \(e^{-u\delta_{\min}}\lesssim\varepsilon\), i.e., \(u\gtrsim \delta_{\min}^{-1}\log(1/\varepsilon)\).
If \(u\) can only reach \(M\Delta u\), we obtain the stated lower bound. This completes the proof.
\end{proof}

\endinput

\paragraph{Support Guardrails and Hierarchical Spectral Localization of Finite Registers: Certificating Gershgorin, Sturm Interlacing, and Schur Residue via Stieltjes}
Following the notation of Proposition \ref{con:xi-shadow-spectrum-determinant-ratio} and Proposition \ref{con:xi-terminating-jfraction-uniqueness},
let \(J\in\RR^{\kappa\times\kappa}\) be the canonical Jacobi certificate of the depth spectrum (symmetric tridiagonal with strictly positive off-diagonal elements), and denote
\[
\beta_i:=J_{ii},
\qquad
\alpha_i:=J_{i,i+1}>0,
\qquad (1\le i\le \kappa-1),
\]
with boundary convention \(\alpha_0=\alpha_\kappa:=0\).
Its spectrum \(\mathrm{Spec}(J)=\{a_j\}_{j=1}^{\kappa}\subset(1,\infty)\) (counting multiplicity) corresponds to depths \(a_j=1+\delta_j\).

\begin{proposition}[Gershgorin Support Guardrail of Jacobi Certificate: Auditable Interval for \texorpdfstring{$(\delta_{\min},\delta_{\max})$}{(delta_min,delta_max)} from Finite Register Coefficients Only]\label{prop:xi-jacobi-gershgorin-support-guardrail}
Let \(a_{\min}:=\min_j a_j\), \(a_{\max}:=\max_j a_j\).
Then there exist explicit support bounds
\[
\boxed{\;
\min_{1\le i\le \kappa}\bigl(\beta_i-\alpha_{i-1}-\alpha_i\bigr)
\ \le\ a_{\min}\ \le\ a_{\max}\ \le\
\max_{1\le i\le \kappa}\bigl(\beta_i+\alpha_{i-1}+\alpha_i\bigr)\; }.
\]
Hence
\[
\boxed{\;
\delta_{\min},\delta_{\max}\ \in\
\Bigl[\min_{i}(\beta_i-\alpha_{i-1}-\alpha_i)-1,\ \max_{i}(\beta_i+\alpha_{i-1}+\alpha_i)-1\Bigr]\; }.
\]
This bound does not depend on any pole localization or full spectral inversion procedure, but only on local coefficients of finite registers, thus serving as an immediately auditable guardrail for in-regime "depth support".
\end{proposition}
\begin{proof}
By Gershgorin circle theorem, any eigenvalue \(\lambda\in\mathrm{Spec}(J)\) must fall in some row's Gershgorin interval
\(
[\beta_i-(|\alpha_{i-1}|+|\alpha_i|),\ \beta_i+(|\alpha_{i-1}|+|\alpha_i|)]
\).
Since \(\alpha_i>0\) and \(\alpha_0=\alpha_\kappa=0\), taking the union yields the stated inequality chain. \qed
\end{proof}

\begin{proposition}[Monotone Contraction of Spectral Localization: First \texorpdfstring{$N$}{N} Register Levels Give Nested Squeezing Intervals for Each Depth Atom \texorpdfstring{$a_j$}{a_j}]\label{prop:xi-jacobi-sturm-nested-localization}
For \(1\le N\le \kappa\) let \(J_N\) be the leading \(N\times N\) principal submatrix of \(J\), and denote its ascending spectrum by
\[
\lambda^{(N)}_1<\cdots<\lambda^{(N)}_N.
\]
Then due to the Sturm interlacing property of tridiagonal Jacobi structure, for each \(1\le N\le \kappa-1\) there is strict interlacing
\[
\boxed{\;
\lambda^{(N+1)}_1<\lambda^{(N)}_1<\lambda^{(N+1)}_2<\cdots<
\lambda^{(N)}_N<\lambda^{(N+1)}_{N+1}\; }.
\]
Therefore, there exists a family of interval indices \(i_N(j)\in\{0,1,\dots,N\}\) (with convention \(\lambda^{(N)}_0:=-\infty\), \(\lambda^{(N)}_{N+1}:=+\infty\)) such that for each \(j\in\{1,\dots,\kappa\}\) and all sufficiently large \(N\)
\[
\lambda^{(N)}_{i_N(j)}\ \le\ a_j\ \le\ \lambda^{(N)}_{i_N(j)+1},
\]
and this squeezing interval contracts hierarchically as \(N\uparrow \kappa\), closing precisely to the point interval \([a_j,a_j]\) at \(N=\kappa\).
Thus "finite certificate length" not only determines defect existence, but also gives auditable localization intervals for each depth atom in hierarchical sense, whose contraction is enforced by tridiagonal spectral theory without relying on numerical heuristics.
\end{proposition}
\begin{proof}[Proof outline]
Strict interlacing follows from three-term recursion of Jacobi matrix characteristic polynomials and Sturm sequence properties.
Interlacing means: each eigenvalue \(\lambda^{(N+1)}_j\) falls in the unique interval \((\lambda^{(N)}_{j-1},\lambda^{(N)}_{j})\);
iterating this localization and tracking the same spectral branch yields hierarchical nested squeezing for the full spectrum \(a_j=\lambda^{(\kappa)}_j\).
When \(N=\kappa\), \(J_N=J\), and the squeezing interval degenerates to a point. \qed
\end{proof}

\begin{proposition}[Positivity of Schur Complement: Stieltjes Type of Finite Register Truncation Error and Full-Scale Sign Certificate]\label{prop:xi-jacobi-truncation-remainder-stieltjes}
Denote
\[
S(t):=\Phi\,\langle e_1,(tI+J)^{-1}e_1\rangle,
\qquad
S_N(t):=\Phi\,\langle e_1,(tI+J_N)^{-1}e_1\rangle,
\qquad (1\le N<\kappa,\ t>0),
\]
where \(\Phi>0\) is the flux constant (consistent with Proposition \ref{con:xi-terminating-jfraction-uniqueness}).
Then there exists a unique finite positive measure \(\nu_N\ge 0\) (supported on \([a_{\min},a_{\max}]\)) such that for all \(t>0\)
\[
\boxed{\;
(-1)^{N-1}\bigl(S(t)-S_N(t)\bigr)
\ =\ \int_{[a_{\min},a_{\max}]}\frac{1}{t+a}\,\dd\nu_N(a)\; }.
\]
Therefore the truncation error is of Stieltjes type at all scales and completely monotone, with sign enforced by parity of \(N\);
and its total mass satisfies
\[
\boxed{\;
\nu_N([a_{\min},a_{\max}])
=\lim_{t\to\infty}t\,(-1)^{N-1}\bigl(S(t)-S_N(t)\bigr)\;}
\]
thus giving an auditable scalar certificate for "unrecoverable residual strength".
\end{proposition}
\begin{proof}[Proof outline]
By Proposition \ref{con:xi-terminating-jfraction-uniqueness}, \(S\) has a terminating Stieltjes continued fraction representation; its \(N\)-th level truncation \(S_N\) is precisely the corresponding convergent.
By Stieltjes continued fraction theory (i.e., Markov--Stieltjes extremal principle and alternating upper/lower bounds of even/odd convergents, see Proposition \ref{prop:xi-markov-stieltjes-extremal-principle}), even/odd convergents bound \(S\) from above/below respectively, and the signed error \((-1)^{N-1}(S-S_N)\) remains a Stieltjes-class function.
Therefore it has the stated positive measure representation; total mass formula follows from Stieltjes asymptotic \(\int (t+a)^{-1}\dd\nu\sim \nu([1,\infty))/t\) as \(t\to\infty\). \qed
\end{proof}

\begin{proposition}[Stieltjes Readout \texorpdfstring{$W_1$}{W1} Stability: Strongest Scale Lipschitz Constant for Depth Measure]\label{prop:xi-stieltjes-w1-lipschitz}
Let \(\mu,\tilde\mu\) be finite positive measures supported on \([1,\infty)\), with equal total mass
\(\mu([1,\infty))=\tilde\mu([1,\infty))\).
Define corresponding Stieltjes orbits
\[
S_\mu(t)=\int\frac{1}{t+a}\,\dd\mu(a),
\qquad
S_{\tilde\mu}(t)=\int\frac{1}{t+a}\,\dd\tilde\mu(a).
\]
Then for any \(t>0\) there is an optimal scale Lipschitz estimate
\[
\boxed{\;
|S_\mu(t)-S_{\tilde\mu}(t)|
\ \le\ \frac{1}{(t+1)^2}\,W_1(\mu,\tilde\mu)\; }.
\]
where \(W_1\) is the first-order Wasserstein distance with cost \(|a-\tilde a|\) (isomorphic definition in Appendix \S\ref{app:kronecker-w1-transport}).
The constant \((t+1)^{-2}\) is given by the strongest Lipschitz constant of kernel \(a\mapsto (t+a)^{-1}\) on \([1,\infty)\), hence scale-optimal.
\end{proposition}
\begin{proof}
Function \(f_t(a):=(t+a)^{-1}\) is differentiable on \([1,\infty)\) and satisfies
\(
|f_t'(a)|=(t+a)^{-2}\le (t+1)^{-2}
\),
hence \(\mathrm{Lip}(f_t)\le (t+1)^{-2}\).
By Kantorovich--Rubinstein duality,
\(
|\int f_t\,\dd(\mu-\tilde\mu)|\le \mathrm{Lip}(f_t)\,W_1(\mu,\tilde\mu)
\),
yielding the conclusion. \qed
\end{proof}

\begin{remark}[Reverse Stability Regime Boundary: Reverse Control of \texorpdfstring{$W_1$}{W1} by Finite Scale Window on Separated Compact Domain]\label{rem:xi-stieltjes-w1-local-inverse}
The above Lipschitz estimate gives the strongest decay law from \(W_1(\mu,\tilde\mu)\) to single-scale readout difference \(|S_\mu(t)-S_{\tilde\mu}(t)|\).
In reverse direction, without imposing a priori compactness and separation constraints, one cannot expect regime-uniform control of transport distance by finite scalar orbit (degeneracy mechanism isomorphic to "depth collision/escape \(\Rightarrow\) inversion ill-posedness").
\par
On parameter domains with finite defect (\(\kappa\) atoms) satisfying the following compactness assumption: support upper bound \(a_{\max}\le A\), minimal separation \(\min_{j\ne k}|a_j-a_k|\ge s\), weight bounds \(0<w_{\min}\le w_j\le w_{\max}\), the map
\(
\mu\mapsto (S_\mu(t))_{t\in[0,T]}
\)
is a finite-dimensional analytic embedding and locally bi-Lipschitz on this compact domain; hence there exists constant \(C=C(\kappa,A,s,w_{\min},w_{\max},T)>0\) such that
\[
W_1(\mu,\tilde\mu)\ \le\ C\int_{0}^{T}(t+1)^2\,\bigl|S_\mu(t)-S_{\tilde\mu}(t)\bigr|\,\dd t.
\]
This inequality rigorously restricts "orbit error \(\Rightarrow\) transport error" auditability within regime-controllable compact domains: once entering collision/escape boundaries, any reverse control constant must inevitably diverge.
\end{remark}

\begin{proposition}[Cauchy--Stieltjes Total Positivity: Variation Diminishing and Certificate Sign Rigidity of Finite Sampling Kernel]\label{prop:xi-cauchy-stieltjes-total-positivity-variation-diminishing}
Let \(1<a_1<\cdots<a_\kappa\), and take strictly increasing sampling \(0\le t_0<\cdots<t_N\).
Define the Cauchy--Stieltjes kernel matrix
\[
\mathcal{C}_{ij}:=\frac{1}{t_i+a_j}\qquad(0\le i\le N,\ 1\le j\le\kappa).
\]
Then \(\mathcal{C}\) is totally positive: all same-order minors are strictly positive.
Furthermore, let \(\mathrm{var}(\cdot)\) denote sign-change count after removing zeros (number of sign changes in sequence), then for any \(x\in\RR^\kappa\) there is variation diminishing property
\[
\boxed{\;
\mathrm{var}\!\bigl(\mathcal{C}x\bigr)\ \le\ \mathrm{var}(x)\; }.
\]
Therefore any finite sampling vector generated by depth spectrum Stieltjes orbit
\(
S(t)=\sum_{j=1}^\kappa \frac{w_j}{t+a_j}
\)
(i.e., \(\mathcal{C}w\)) satisfies an in-regime sign rigidity: when performing any linear residual/test on it, the number of visible sign oscillations is strictly bounded by the sign structure and dimension \(\kappa\) of input coefficients, thus excluding "high-frequency pseudo-oscillation" certificate artifacts.
\end{proposition}
\begin{proof}[Proof outline]
Take any same-order index subsets \(I=\{i_1<\cdots<i_m\}\subseteq\{0,\dots,N\}\), \(J=\{j_1<\cdots<j_m\}\subseteq\{1,\dots,\kappa\}\).
The corresponding \(m\times m\) minor is the classical Cauchy determinant:
\[
\det\!\Bigl(\frac{1}{t_{i_p}+a_{j_q}}\Bigr)_{1\le p,q\le m}
=
\frac{\prod_{1\le p<q\le m}(t_{i_q}-t_{i_p})\cdot\prod_{1\le p<q\le m}(a_{j_q}-a_{j_p})}
{\prod_{p,q}(t_{i_p}+a_{j_q})}.
\]
Since \(t_{i_q}-t_{i_p}>0\), \(a_{j_q}-a_{j_p}>0\) and \(t_{i_p}+a_{j_q}>0\), the minor is strictly positive, thus \(\mathcal{C}\) is totally positive.
Variation diminishing property is a direct consequence of the fundamental theorem for totally positive matrices (Gantmacher--Krein variation diminishing theorem). \qed
\end{proof}

\endinput


\begin{proposition}[Minimal Direct Product Theorem for Depth--Height: Covariance Forces Amplitude/Phase Separation]
In finite defect caliber, if a family of auditable spectral fingerprints simultaneously satisfies:
\begin{enumerate}
  \item Additive over defects;
  \item Strong continuous covariance under height translation \(\gamma\mapsto\gamma+x_0\);
  \item Covariance under depth-side two-parameter semi-direct action \((\tau,t)\) (Proposition (Unique Commutation Law of Two-Parameter Renormalization Semigroup));
\end{enumerate}
then its structure must decompose as "depth amplitude part \(\times\) height phase part" direct product representation: amplitude enters only through Stieltjes/Laplace orbit of depth measure \(\mu\) (equivalently \(S\) or \(Z\)), while height can only enter through translation-covariant phase carrier \(\sigma\)'s character \(e^{-\mathrm{i}\gamma\xi}\) (or its discretization).
In particular, when retaining only scale entropy curve \(S\) (or equivalent \(Z\)), height parameter becomes indistinguishable in regime sense; any construction attempting to endogenously recover height side from depth side will break one of the above three structural constraints, thus degenerating to recycling (congruence folding or \(\mathsf{NULL}\)) in auditable closure.
\end{proposition}

\endinput


\endinput
